
\documentclass[UTF8,a4paper,12pt]{ctexart}

\usepackage{geometry}
\usepackage{amsmath,amssymb}
\usepackage{booktabs,longtable,array}
\usepackage{graphicx,float}
\usepackage{xcolor}
\usepackage{enumitem}
\usepackage{caption}
\usepackage{subcaption}
\usepackage[hidelinks]{hyperref}

\geometry{left=2.6cm,right=2.6cm,top=2.5cm,bottom=2.5cm}
\setlength{\parindent}{2em}
\setlength{\parskip}{0.25em}
\renewcommand{\arraystretch}{1.18}
\setcounter{tocdepth}{2}
\setcounter{secnumdepth}{3}

\newcommand{\figplaceholder}[2][5.2cm]{%
  \begin{figure}[H]
    \centering
    \fbox{\parbox[c][#1][c]{0.92\textwidth}{\centering #2}}
  \end{figure}
}

\title{山区洪涝救援中无人机运输--通信协同调度与独立资源配置\\
\large 当前整合草稿（第1--9章）}
\author{}
\date{基于 main@95c2b4a24409fbce6c216b45ed02c69a7612634c 的冻结结果口径}

\begin{document}

\maketitle


\begin{abstract}
针对山区洪涝灾害中道路受阻、地形起伏、通信受限及救援资源异构等条件下的无人机应急运输问题，本文研究从单架次安全运输到多任务调度、连续通信保障及独立任务组资源配置的逐层优化。不同于将四个问题分别处理的做法，本文首先建立统一的“地形--载荷--时间--能耗”运输任务模型，利用AEQD坐标和DEM完整像元遍历确定航段距离与巡航高度，并结合逐段载荷变化、非线性载荷--航程关系和返航安全余量计算任务时间与能耗；随后依次加入不可拆货箱组批、实体无人机--共享电池时序、中继通信以及独立组织边界等约束，从而在同一物理口径下分析系统可执行性。

对于单点运输问题，在20\%返航安全余量下计算45组“服务区--机型”安全载荷，并通过精确组批得到最少运输架次数为18，达到严格下界；在此条件下总能耗为59.131296 kWh，累计作业时间为546.266929 min。进一步的安全余量分析发现，连续变化的安全要求会通过不可拆货箱结构触发离散架次数跃迁，说明额定载荷不能直接代表山区任务中的实际运输能力。对于多架次运输调度问题，在有限实体无人机和共享电池条件下构造零加权迟到的23架次方案，A/B/C型任务数为11/6/6，最后返航时间为5693.231489 s，总能耗为66.214460 kWh。通过全局有效下界将最优完成时间限制在
\[
5433.956428\le T^*\le5693.231490\ {\rm s},
\]
以上界为分母的认证间隙为4.5541\%，因此该方案具有量化近优性证据，但不宣称完成时间已达到全局最优。另构造的22架次可行见证工期为7767.058323 s，表明减少架次数并不等价于缩短系统完成时间。

进一步考虑山区连续通信后，将实际运输轨迹上的直连失效区间转化为中继时空服务需求，并联合调度中继位置、服务时段、中继无人机及能源组件。正式方案由23个运输架次和4个中继架次组成，联合最后返航时间为6340.439339 s，总能耗为69.308531 kWh。在固定运输轨迹与5755个候选中继位置构成的有限域内，利用196个真实通信需求的覆盖证书证明至少需要4个中继架次；0.50 dB和0.55 dB额外统一传播损耗下仍存在完整可行方案，量化了通信鲁棒性与任务工期之间的代价。关键任务调整还使电池提前57.025 s恢复并使联合工期缩短约54.652 s，说明系统工期更直接受关键资源恢复链影响，而非仅由总能耗决定。

在冻结上述运输与通信计划后，进一步研究独立任务组的资源配置。严格模型下，服务区因运输和中继关联形成3个不可拆任务块；两组和三组合法分区分别只有3种和1种，可进行完整枚举。两组最小资源缺口为2，即增加1架C型运输无人机和1组C型电池；三组最小缺口为7。值得注意的是，两组总配置需求仅为29，小于原始库存总数30，但由于资源类型不匹配仍无法实现独立执行，表明异构系统中资源总量不能替代分类型资源结构判断。综合四问可见，山区无人机救援中的单架次安全、系统工期、通信连续性和独立资源充足性并不具有自动传递关系；真正决定系统性能的是运输、能源、通信与组织资源在统一时间轴上的恢复、复用与协同关系。本文结果可为应急任务中的安全余量设定、关键资源识别、中继鲁棒配置及分类型装备补充提供量化依据。
\end{abstract}

\noindent\textbf{关键词：}山区洪涝救援；无人机应急运输；异构资源调度；连续通信保障；资源池化；多目标优化


\begin{center}
\small
本文件为当前已完成内容的单文件 \LaTeX{} 整合稿，包含第1--9章。；图均保留占位，正式图后续替换。
Q2正式证据采用 Q2-EVIDENCE-V7-FOCUSED：$5433.956428\le T^*\le5693.231490$ s，
UB归一化间隙为 $4.5541\%$。
\end{center}

\tableofcontents
\newpage


% ============================================================
\section{问题分析}
\label{sec:problem_analysis}

\subsection{问题背景与任务本质}

山区洪涝灾害发生后，道路中断、地形复杂以及通信基础设施受损，会使传统地面运输难以及时覆盖受灾区域。无人机具有部署速度快、对道路依赖低和可快速穿越复杂地形等特点，因此可以承担应急物资运输和临时通信保障任务\cite{dorling2017,altheeb2017relief}。

但在实际救援过程中，“无人机能飞到”并不意味着整个救援任务就已经可执行。对于给定的运输基地、15个受灾服务区、80个不可拆货箱以及多种异构无人机，任务同时受到物理能力、组合决策、资源时序、通信保障与组织独立等多类约束。

因此，本题并不是单纯的无人机路径规划、装箱或调度问题，而是一个由运输能力、任务组合、资源调度、通信保障和组织划分逐层耦合形成的多阶段优化问题。四个问题不是平行的四道小题，而是在同一物理系统上逐层加入新的现实约束：
\[
\boxed{
\text{单架次安全能力}
\rightarrow
\text{系统运输调度}
\rightarrow
\text{运输--通信联合执行}
\rightarrow
\text{独立组织资源配置}
}.
\]

\subsection{四个问题的递进关系}

\subsubsection{问题一：先回答“一趟到底能运什么”}

问题一中，每一个运输架次只服务一个受灾区域，不考虑实体无人机之间的并行执行，也不考虑共享电池库存与充电。因此核心首先是确定：在给定服务区地形和机型参数后，一架无人机究竟能够安全携带多少货物。

无人机额定载荷不能直接作为山区条件下的实际载荷上限，因为远距离航程、巡航高度和爬升过程都会增加能耗。当返航安全余量被考虑后，不同“机型--服务区”组合具有不同最大安全载荷。

进一步地，实际物资由质量、体积不同且不可拆分的货箱构成，因此还需要完成
\[
\boxed{
\text{连续运输能力}
\rightarrow
\text{离散货箱组批}
}.
\]
问题一的本质可概括为“能力边界确定+不可拆货箱精确组批”。

\subsubsection{问题二：单架次可行不等于整个机群能及时完成}

问题二开始考虑真实有限资源：不同类型实体无人机数量有限，同类型任务共享有限电池，电池返航后需要继续充电，不同货箱还具有送达时限。因此
\[
\boxed{
\text{单架次可行}
\not\Rightarrow
\text{整个运输计划可执行}
}.
\]
问题二由此升级为货箱组批、多点路线、异构机型、实体飞机、共享电池和任务时限共同作用的联合调度问题，其目标也由“尽量减少架次”转变为“首先满足时限，再尽可能缩短整个机群最后返航时间”。

\subsubsection{问题三：运输可执行不等于全过程通信可用}

问题二仍隐含“运输无人机执行过程中始终能够与地面保持通信”的条件。山区地形可能使无人机在飞行中途发生通信遮挡，因此
\[
\boxed{
\text{运输排程可执行}
\not\Rightarrow
\text{运输全过程通信可用}
}.
\]
当直连可用时优先直连；当直连失效时，则需要中继无人机在正确的位置和正确时间提供服务。中继无人机本身又受到位置、到达时间、服务窗口、实体数量和能源组件约束，因此问题三本质上是运输与通信的时空联合调度。

\subsubsection{问题四：全局共享可执行不等于独立分组后资源仍充足}

前三问均默认全部资源属于统一资源池。问题四进一步要求将服务区划分为若干相互独立的任务组，一旦不同组之间不允许共享资源，第三问中原有的跨服务区资源复用链就可能被切断。因此
\[
\boxed{
\text{全局资源充足}
\not\Rightarrow
\text{各独立任务组资源充足}
}.
\]
问题四真正研究的是组织边界改变以后资源池化收益损失多少，以及为了保证各组独立执行需要额外增加哪些类型资源。

\subsection{全文核心矛盾}

综合四问，本题最核心的结构是不同层次的“最优”并不具有传递性：
\[
\text{单次最省电}
\not\Rightarrow
\text{整个任务最快完成},
\]
\[
\text{架次数最少}
\not\Rightarrow
\text{总能耗最低},
\]
\[
\text{静态通信覆盖良好}
\not\Rightarrow
\text{动态中继方案可执行},
\]
\[
\text{资源总量充足}
\not\Rightarrow
\text{独立分组后每种资源都足够}.
\]
因此本文不构造一个将架次数、时间、能源、通信裕度和资源均衡强行合并的单一加权目标，而是根据不同问题所处的决策层次分别建立清晰的目标优先关系，并通过Pareto方案、鲁棒备选和敏感性分析展示实际权衡。

\subsection{总体建模与证据思路}

为避免四问分别建立割裂模型，本文首先构造统一底层运输任务物理模型，包括AEQD航段、DEM完整像元遍历、巡航高度、三阶段飞行、逐段剩余载荷、载荷相关航程、水平与爬升能耗、返航安全余量以及完整作业时间。

在此基础上，Q1增加安全载荷和不可拆货箱组批；Q2增加实体无人机、共享电池、充电恢复和送达时限；Q3增加连续通信、中继位置和中继资源；Q4冻结Q3正式任务，仅改变服务区组织划分并重新计算独立资源需求。

不同问题采用与其结构匹配的证据方式：Q1使用精确MILP和DP；Q2同时构造完整可行上界与全局有效下界；Q3采用连续通信验证与有限候选域下界；Q4在任务块压缩后直接完整枚举。全文严格区分“全局最优”“有效上下界”“有限域最优”“完整可行见证”和“局部未发现改善”等不同证据等级。

\figplaceholder{图1-1占位：四个问题的约束递进关系与统一建模框架。}

% ============================================================
\section{模型假设、符号与评价口径}
\label{sec:assumptions_symbols}

\subsection{条件分类与建模口径}

为避免将题目给定条件、建模假设和执行策略混淆，本文将条件分为四类。

\paragraph{题面给定条件}
包括服务区和基地坐标、地面高程、DEM数据、80个货箱的质量体积和目标服务区、A/B/C三类运输无人机参数、实体无人机与电池库存、载荷--航程关系、飞行速度、准备装卸交接时间、充电规则、通信链路阈值以及中继资源参数。

\paragraph{统一建模口径}
包括以O01为中心建立WGS84 AEQD米制平面、水平航迹按AEQD直线处理、巡航高度取航段经过DEM像元最高值加50 m、服务区作业高度为地面高程加30 m、每次离开服务区后重新爬升至下一航段巡航高度，以及所有底层数值内部使用完整精度。

\paragraph{执行策略}
包括任务准备前占用对应实体无人机和兼容满电电池；无人机占用持续至返航；电池占用持续至返航后恢复满电；同型资源共享、不同类型不可互换；中继无人机返航后计入规定周转时间；中继能源组件持续占用至恢复；最终充电不计入当前任务makespan但仍计入资源占用。

\paragraph{反事实与敏感性情景}
包括Q1统一改变返航安全余量、Q3增加0.50/0.55 dB统一额外通信损耗、Q4允许跨组完整复制中继任务以及Q4某类初始库存减少1个。此类结果均明确标为敏感性或反事实，不替代正式主答案。

\subsection{基本模型假设}

\begin{enumerate}[label=\textbf{假设\arabic*.},leftmargin=3.5em]
\item 货箱不可拆分，每个原始货箱必须作为完整单位运输。
\item 无人机货物约束包括总质量和总体积，不进一步考虑长宽高方向的三维几何装箱。
\item 任意两个访问节点之间采用统一AEQD平面直线飞行，不额外规划题目未给出的避障绕行、建筑、树木、禁飞区或动态障碍物。
\item 节点地面高程采用题目原始值，DEM主要用于计算航段飞越地形最高值。
\item 下降阶段计入飞行时间，但不额外增加下降附加能耗；返程即使空载仍计算水平和重新爬升能耗。
\item 主模型不额外引入风场、降雨阻力、温度、电池老化或随机设备故障，因此结论属于题面给定模型下的确定性结果。
\item 同型号无人机、同型号电池和同类中继资源视为性能一致，可在同类型内部替换。
\item 不同资源类型不能互换，例如A型与C型运输机、C型电池与中继能源组件不能相互替代。
\item 多组电池可并行充电，不额外设置共享充电器容量约束。
\item 通信采用直连优先、单中继补偿，不构造两级以上多跳中继网络。
\item Q4正式模型严格冻结Q3 \texttt{time} 方案的箱组、机型、路线、绝对时间、中继位置、服务窗口和实际保障关系；中继复制仅作为敏感性分析。
\end{enumerate}

\subsection{主要符号说明}

全文仅在本节列跨章节反复使用的核心符号，局部算法变量在对应章节定义。

\begin{longtable}{p{3cm}p{8.5cm}p{2cm}}
\caption{主要集合、任务与物理符号}\\
\toprule
符号 & 含义 & 单位\\
\midrule
\endfirsthead
\toprule
符号 & 含义 & 单位\\
\midrule
\endhead
$\mathcal V$ & 基地与服务区节点集合 & --\\
$\mathcal S$ & 15个服务区集合 & --\\
$\mathcal B$ & 全部货箱集合 & --\\
$\mathcal B_i$ & 服务区$i$的货箱集合 & --\\
$\mathcal M$ & 运输无人机机型集合$\{A,B,C\}$ & --\\
$\mathcal R$ & 运输任务集合 & --\\
$\mathcal U_m$ & 机型$m$的实体运输无人机集合 & --\\
$\mathcal K_m$ & 机型$m$的运输电池集合 & --\\
$\mathcal U^R$ & 中继无人机集合 & --\\
$\mathcal K^R$ & 中继能源组件集合 & --\\
$\mathcal P^R$ & 中继悬停候选位置集合 & --\\
$\mu_b$ & 货箱$b$质量 & kg\\
$\nu_b$ & 货箱$b$体积 & m$^3$\\
$Q_m$ & 机型$m$额定载荷 & kg\\
$V_m$ & 机型$m$货舱体积 & m$^3$\\
$E_m$ & 机型$m$可用电池能量 & kWh\\
$L_m(q)$ & 载荷$q$下等效航程 & m\\
$\alpha_m$ & 返航安全余量 & --\\
$d_{ij}$ & 节点$i,j$水平距离 & m\\
$H_{ij}$ & 航段巡航海拔 & m\\
$r$ & 一项完整运输任务 & --\\
$m_r$ & 任务$r$使用机型 & --\\
$\mathcal B_r$ & 任务$r$携带货箱集合 & --\\
$\pi_r$ & 任务$r$有序访问路径 & --\\
$q_{r,j}$ & 第$j$航段实际载荷 & kg\\
$p_r$ & 任务完整作业时长 & s\\
$E_r$ & 任务总能耗 & kWh\\
$SOC_r^{return}$ & 任务返航SOC & --\\
$\delta_{rb}$ & 箱$b$相对任务开始的送达偏移 & s\\
$s_r$ & 任务开始时刻 & s\\
$C_b$ & 箱$b$实际送达时刻 & s\\
$M_{ab}(t)$ & 链路$a-b$在时刻$t$的通信裕度 & dB\\
$\Delta$ & 统一额外传播损耗 & dB\\
$I_r^U$ & 运输无人机占用区间 & s\\
$I_r^B$ & 运输电池占用区间 & s\\
$I_q^R$ & 中继无人机占用区间 & s\\
$I_q^K$ & 中继能源组件占用区间 & s\\
$N_{gp}^{min}$ & 组$g$对资源类型$p$的最小需求 & 个\\
$I_p$ & 类型$p$原始库存 & 个\\
$D_p$ & 分组后类型$p$总需求 & 个\\
$G_p$ & 类型$p$资源缺口 & 个\\
\bottomrule
\end{longtable}

\subsection{时间区间与数值约定}

所有实体资源占用统一采用半开区间
\[
[start,release).
\]
若前一任务恰好在时刻$t$释放资源，而后一任务在同一时刻$t$开始，则二者不被错误判断为冲突。

内部计算统一使用：
\[
\text{距离：m},\quad
\text{时间：s},\quad
\text{质量：kg},\quad
\text{体积：m^3},\quad
\text{能量：kWh},\quad
\text{通信损耗：dB}.
\]
正文中的分钟和小数位仅用于展示，后续计算始终继承未舍入数值。因此
\[
\boxed{
\text{显示精度降低}
\neq
\text{计算精度降低}
}.
\]

\subsection{各问题评价目标}

Q1采用
\[
\boxed{
\min_{\rm lex}(N,E,T^{sum})
}
\]
即最少运输架次、固定架次后的最低能耗和最短累计作业时间，并另外计算Pareto目标向量。

Q2采用
\[
\boxed{
\min_{\rm lex}(L_w,T_{\max},E,N)
}
\]
即先最小加权迟到，再最小系统完成时间、能耗和架次数。由于已有完整可行解达到$L_w=0$且$L_w\ge0$，第一层全局最优值为
\[
L_w^*=0.
\]

Q3首先继续保持$L_w=0$和硬截止可行，正式 \texttt{time} 方案主要最小化联合完成时间 $T_{\max}^{joint}$。通信裕度$\Delta$不与工期和能耗强行加权，而通过$\Delta=0,0.50,0.55$等独立情景比较。

Q4正式严格模型采用
\[
\boxed{
\text{资源缺口}
\rightarrow
\text{总配置}
\rightarrow
\text{工作量CV}
\rightarrow
\text{规范ID}
}
\]
的词典序选择，避免将设备数量、工作量和异构缺口强行折算成单一综合分值。

\subsection{可行性与最优性表述规范}

若方案满足当前问题全部约束，只称为“完整可行方案”，并不自动说明全局最优。

只有当存在严格下界并有可行解达到该下界，或精确算法完整覆盖原可行域并证明最优时，才使用“全局最优”表述。例如Q1最少架次满足
\[
LB=UB=18.
\]

若只得到
\[
LB\le T^*\le UB,\qquad LB<UB,
\]
则称为“具有认证上下界的当前最优已知可行解”，Q2属于该情形。

若证明只针对明确有限候选集合成立，则写成“给定候选域中的最优或下界”，不能自动推广到连续自由空间。Q3的4中继架次下界仅对应固定运输轨迹和5755个候选位置。

若有限邻域、有限时间或启发式搜索没有找到更优方案，只表述为“在给定搜索范围内未发现改善”。求解器达到时间限制时统一记为unknown，不将其解释为不可行。

\subsection{验证、对照与敏感性使用原则}

本文将三类实验严格区分。

“验证”回答结果是否算对，例如独立物理重算、MILP/DP交叉验证、通信端点检查和资源峰值多方法一致性。

“对照实验”回答当前模型或方法相对合理基线多解决了什么，例如单机型与混合机型、22/23架次、静态通信评分与完整联合排程。

“敏感性分析”回答关键参数变化后决策是否发生结构变化，例如安全余量、通信额外损耗、初始库存以及是否允许复制中继任务。

因此，三类证据不统一笼统称为“模型验证”。

\subsection{本章小结}

本章区分了题面条件、统一建模口径、执行策略和反事实敏感性，并统一了全文符号、单位、资源占用区间和评价目标。根据不同问题所处决策层次分别设置目标优先级，同时根据证据强度区分全局最优、有效上下界、有限域最优、完整可行和局部未改善，为第3章统一物理模型以及第4--7章逐层优化提供共同口径。


% ============================================================
\section{数据处理与统一运输任务物理模型}
\label{sec:foundation}

四个问题虽然研究的决策层次不同，但底层均建立在同一套救援场景、货箱数据、地形数据和无人机物理参数之上。若分别在各问题中重复建立距离、航段高度、载荷、飞行时间和能耗模型，不仅造成正文重复，还可能导致不同问题采用不一致的物理口径。因此，本文首先建立统一的运输任务物理模型。

对于后续任意运输任务，均采用如下计算链：
\[
\boxed{
\text{节点与货箱}
\rightarrow
\text{空间航段}
\rightarrow
\text{DEM地形}
\rightarrow
\text{逐段载荷}
\rightarrow
\text{飞行时间}
\rightarrow
\text{能耗}
\rightarrow
\text{任务属性}
}
\]
在此基础上，不同问题只增加各自特有的决策层：问题一增加单点安全载荷和不可拆货箱组批；问题二增加实体无人机、共享电池和配送时限；问题三进一步增加连续通信和中继任务；问题四则冻结问题三的完整任务计划，分析独立分组后的资源配置。

\subsection{基础数据对象}

记运输基地为 O01，15个救援服务区集合为
\[
\mathcal S=\{S001,S002,\ldots,S015\},
\]
统一记全部空间节点集合为
\[
\mathcal V=\{O01\}\cup\mathcal S.
\]
每个节点 $i\in\mathcal V$ 具有原始经纬度坐标 $(\lambda_i,\varphi_i)$ 和题目给定地面高程 $z_i$。

全部救灾物资由80个不可拆货箱组成，记货箱集合为 $\mathcal B$。对于任意货箱 $b\in\mathcal B$，定义 $\mu_b$ 为质量，$\nu_b$ 为体积，$s(b)$ 为目标服务区。全部货箱总质量为 $758$ kg，总体积为 $2.011$ m$^3$。

运输无人机类型集合记为
\[
\mathcal M=\{A,B,C\}.
\]
三类机型的主要参数如下。

\begin{table}[H]
\centering
\caption{三类运输无人机主要参数}
\label{tab:common_drone}
\begin{tabular}{lrrr}
\toprule
参数 & A型 & B型 & C型\\
\midrule
含电池空载质量/kg & 70 & 65 & 69.9\\
额定载荷/kg & 25 & 30 & 80\\
货舱体积/m$^3$ & 0.060 & 0.073 & 0.250\\
巡航速度/(m/s) & 12 & 15 & 15\\
空载航程/m & 25000 & 28000 & 26000\\
满载航程/m & 20000 & 16000 & 12000\\
可用电量/kWh & 4.5 & 4.0 & 8.0\\
基准返航余量 & 20\% & 20\% & 20\%\\
\bottomrule
\end{tabular}
\end{table}

实体无人机数量和共享电池数量从问题二开始进入资源排程，因此本章只负责描述“一项运输任务需要多少时间和能源”，动态资源占用模型在第5章展开。

\subsection{统一空间坐标与DEM航段模型}

以运输基地O01作为中心建立WGS84方位等距投影（AEQD）\cite{snyder1987}：
\[
(x_i,y_i)=\operatorname{AEQD}_{O01}(\lambda_i,\varphi_i).
\]
任意两个节点 $i,j$ 之间的水平飞行路径定义为AEQD平面中的直线
\[
\mathbf r_{ij}(u)
=
(1-u)
\begin{bmatrix}
x_i\\y_i
\end{bmatrix}
+
u
\begin{bmatrix}
x_j\\y_j
\end{bmatrix},
\qquad u\in[0,1],
\]
其水平距离为
\[
d_{ij}
=
\sqrt{(x_j-x_i)^2+(y_j-y_i)^2}.
\]

对于航段 $i\rightarrow j$，将AEQD平面直线连续反投影到原始经纬度坐标系，并与题目给定DEM栅格逐像元求交。记与航段实际相交的DEM像元集合为 $\mathcal C_{ij}$，则航段经过的最大地形高程为
\[
z_{ij}^{\max}
=
\max_{c\in\mathcal C_{ij}}DEM(c).
\]
本文采用完整像元相交而不是稀疏采样或thin Bresenham近似，以避免遗漏短暂经过的高程像元。

航段巡航海拔统一设为
\[
H_{ij}=z_{ij}^{\max}+50.
\]
运输基地O01的作业高度采用其地面高程
\[
h_{O01}=z_{O01},
\]
服务区的物资交接高度设为
\[
h_i=z_i+30,\qquad i\in\mathcal S.
\]
因此
\[
h_{ij}^{\uparrow}=H_{ij}-h_i,\qquad
h_{ij}^{\downarrow}=H_{ij}-h_j.
\]
无人机在某服务区完成交付后若继续飞往下一服务区，必须重新经历爬升、巡航和下降。

\subsection{统一运输任务与逐段载荷演化}

将任意运输任务统一表示为
\[
r=(m_r,\mathcal B_r,\pi_r),
\]
其中 $m_r\in\mathcal M$ 为任务机型，$\mathcal B_r\subseteq\mathcal B$ 为货箱集合，
\[
\pi_r=(v_0,v_1,\ldots,v_K,v_{K+1})
\]
为有序访问路径，且 $v_0=v_{K+1}=O01$。

任务初始载荷与总体积为
\[
q_{r,0}=\sum_{b\in\mathcal B_r}\mu_b,\qquad
V_r=\sum_{b\in\mathcal B_r}\nu_b,
\]
并满足
\[
q_{r,0}\le Q_{m_r},\qquad
V_r\le V_{m_r}.
\]

若在服务区 $v_j$ 交付货物质量 $D_{rj}$，则
\[
q_{r,j}=q_{r,j-1}-D_{rj},
\]
从而
\[
q_{r,0}\ge q_{r,1}\ge\cdots\ge q_{r,K}=0.
\]
航段 $v_j\rightarrow v_{j+1}$ 使用的实际载荷为 $q_{r,j}$。因此访问顺序不仅影响空间距离，也影响载荷演化及能耗。

\subsection{统一飞行与作业时间模型}

任意航段 $i\rightarrow j$ 的飞行时间为
\[
t_{ij}^{fly}(m)
=
\frac{h_{ij}^{\uparrow}}{v_m^{\uparrow}}
+
\frac{d_{ij}}{v_m^c}
+
\frac{h_{ij}^{\downarrow}}{v_m^{\downarrow}}.
\]
任务纯飞行时间为
\[
T_r^{fly}
=
\sum_{j=0}^{K}
t_{v_jv_{j+1}}^{fly}(m_r).
\]

设机型 $m$ 的固定准备时间、单箱装载时间、每站基础交接时间和单箱交接附加时间分别为
$t_m^{prep}$、$t_m^{load}$、$t_m^{hand}$、$t_m^{box}$。
若任务 $r$ 共运输 $n_r$ 个货箱并访问 $K$ 个服务区，第 $j$ 个服务区交付 $n_{rj}$ 个货箱，则任务完整作业时长为
\[
p_r
=
t_{m_r}^{prep}
+
n_rt_{m_r}^{load}
+
T_r^{fly}
+
\sum_{j=1}^{K}
\left(
t_{m_r}^{hand}
+
n_{rj}t_{m_r}^{box}
\right).
\]

沿任务路径累计即可得到每个货箱相对任务开始的交付偏移 $\delta_{rb}$。若任务实际开始时间为 $s_r$，则
\[
C_b=s_r+\delta_{rb}.
\]

\subsection{载荷相关航程与能耗模型}

无人机配送研究表明，载荷、航程与能耗之间存在显著耦合关系\cite{dorling2017,zhang2021energy,zhang2023corrigendum}。本文具体的载荷--航程函数及指数仍严格沿用题面给定关系。对机型 $m$，载荷 $q$ 下的等效航程统一表示为
\[
\boxed{
L_m(q)
=
L_m^0
-
(L_m^0-L_m^F)
\left(\frac{q}{Q_m}\right)^{3/2}
},
\qquad 0\le q\le Q_m.
\]
水平飞行能耗为
\[
E_m^{hor}(d,q)
=
E_m\frac{d}{L_m(q)}.
\]
若爬升高度为 $h$，则爬升附加能耗为
\[
E_m^{up}(q,h)
=
\frac{(m_m^0+q)gh}{\eta_m\,3.6\times10^6},
\qquad g=9.81\ {\rm m/s^2}.
\]
下降阶段仅计飞行时间，不另加附加能耗。

因此任务 $r$ 的第 $j$ 个航段能耗为
\[
E_{rj}
=
E_{m_r}^{hor}
(d_{v_jv_{j+1}},q_{r,j})
+
E_{m_r}^{up}
(q_{r,j},h_{v_jv_{j+1}}^\uparrow),
\]
整任务能耗为
\[
E_r=\sum_{j=0}^{K}E_{rj}.
\]
最后返航航段虽为空载，但仍具有水平巡航能耗和重新爬升能耗，因此
\[
\boxed{\text{空载返程}\neq\text{零返程能耗}}.
\]

\subsection{安全余量与返航SOC}

若机型 $m$ 的返航安全余量为 $\alpha_m$，则任务可使用能量预算为
\[
E_m^{budget}=(1-\alpha_m)E_m.
\]
任务 $r$ 必须满足
\[
E_r\le(1-\alpha_{m_r})E_{m_r}.
\]
返航SOC定义为
\[
SOC_r^{return}
=
1-\frac{E_r}{E_{m_r}},
\]
故安全条件等价于
\[
SOC_r^{return}\ge\alpha_{m_r}.
\]

基准场景统一取 $\alpha_m=20\%$。第4章利用该条件求单点最大安全载荷，第5、6章用于判断多点运输任务的能源可行性，第7章直接继承第6章冻结任务。

\subsection{统一任务输出接口与评价量}

经过上述模型，每个运输任务统一转换为
\[
r
\longrightarrow
\left(
m_r,\mathcal B_r,\pi_r,p_r,E_r,SOC_r^{return},\{\delta_{rb}\}
\right).
\]
后续各章只在这一公共接口上增加新的系统约束。

若方案包含运输任务集合 $\mathcal R$，则
\[
N=|\mathcal R|,
\qquad
E^{trans}=\sum_{r\in\mathcal R}E_r,
\qquad
T^{sum}=\sum_{r\in\mathcal R}p_r.
\]
若任务存在实际开始时刻，则系统最后完成时间定义为
\[
T^{max}
=
\max_{r\in\mathcal R}(s_r+p_r).
\]
因此
\[
\boxed{T^{sum}\neq T^{max}}.
\]

若货箱 $b$ 的期望送达时刻为 $d_b$，迟到量与加权迟到分别为
\[
L_b=\max(0,C_b-d_b),
\qquad
L_w=\sum_{b\in\mathcal B}\omega_bL_b.
\]

\subsection{公共模型验证与数值口径}

公共模型一旦出错会同时影响Q1--Q4，因此正文只保留公共层面的独立核验：AEQD/GeographicLib距离与DEM相交像元采用不同方法交叉检查；载荷相关航程、水平能耗、爬升能耗和返程空载状态使用独立公式重算；所有正式任务统一检查货箱集合、质量、体积、逐段载荷、时间、能耗与SOC。

内部统一使用m、s、kg、m$^3$和kWh。论文中的分钟与小数位只用于展示，后续章节继承未舍入数值。完整原始字段、文件哈希、故障注入、清洁复现和逐项数值审计放入附录，不在各问题章节重复展开。

\subsection{公共模型与四问接口}

\begin{table}[H]
\centering
\caption{统一公共模型与四个问题的接口}
\label{tab:common_interface}
\begin{tabular}{p{1.3cm}p{6.0cm}p{7.0cm}}
\toprule
问题 & 直接继承第3章 & 新增决策层\\
\midrule
Q1 & 单点航段、时间、能耗、安全余量 & 最大安全载荷、不可拆货箱组批、Pareto、余量临界事件\\
Q2 & 多点路径、逐段载荷、任务时长、能耗、箱级送达偏移 & 实体飞机、共享电池、充电恢复、时限、联合调度、上下界\\
Q3 & 完整运输任务物理属性 & 连续通信、中继选址、中继任务与联合排程\\
Q4 & 冻结后的Q3任务时间与资源占用 & 不可拆分组、组内资源峰值、库存缺口、池化损失\\
\bottomrule
\end{tabular}
\end{table}

% ============================================================
\section{问题一：单点安全运输能力与货箱组批优化}
\label{sec:q1}

\subsection{问题结构与数学转化}

问题一研究单点往返条件下的运输能力和离散组批；多架次无人机配送中的路线、载荷与能源耦合也是相关文献长期关注的核心问题\cite{murray2015fstsp,dorling2017}。对任意服务区 $i$ 和机型 $m$，首先由第3章统一物理模型确定最大安全载荷 $w_{im}^{safe}$；随后将连续载荷能力转化为满足质量、体积和能耗约束的离散货箱批型；最后在每个服务区内部进行精确组批优化。由于Q1禁止跨服务区组批，且不存在实体飞机和共享电池之间的跨任务竞争，因此15个服务区可独立求解后汇总。

\figplaceholder{图4-1占位：问题一从物理能力层、可行批型层到精确组批优化层的数学转化框架。}

\subsection{最大安全载荷}

对服务区 $i$、机型 $m$，完整往返能耗记为 $E_{im}(w)$，则
\[
w_{im}^{safe}
=
\max\left\{
w\in[0,Q_m]:
E_{im}(w)\le(1-\alpha_m)E_m
\right\}.
\]
由于载荷增加会降低等效航程并增加爬升能耗，$E_{im}(w)$ 关于 $w$ 单调增加，因此能量约束在额定载荷之前生效时具有唯一边界根。

在20\%返航安全余量下，共得到45个“服务区--机型”安全载荷，其中39项达到额定质量上限，6项首先受到能源约束。典型能源受限值包括：
\[
w_{S002,C}^{safe}\approx68.860,\quad
w_{S003,C}^{safe}\approx68.206,
\]
\[
w_{S004,C}^{safe}\approx63.697,\quad
w_{S008,B}^{safe}\approx28.801,
\]
\[
w_{S008,C}^{safe}\approx58.904,\quad
w_{S012,C}^{safe}\approx68.117\ {\rm kg}.
\]
这说明额定载荷并不能完全代表山区条件下的实际运输能力。

\figplaceholder[6.0cm]{图4-2占位：15个服务区$\times$3种机型的低饱和蓝灰安全载荷矩阵。颜色表示额定能力损失率，单元格显示最大安全载荷/kg。}

\subsection{可行批型与组批模型}

对服务区 $i$ 的任意非空货箱子集 $S\subseteq\mathcal B_i$，若机型 $m$ 同时满足
\[
\sum_{b\in S}\mu_b\le Q_m,
\qquad
\sum_{b\in S}\nu_b\le V_m,
\]
以及
\[
E_{im}\left(\sum_{b\in S}\mu_b\right)
\le(1-\alpha_m)E_m,
\]
则 $(m,S)$ 为可行批型。

定义二元变量 $y_p$ 表示是否选择批型 $p$。每个货箱必须恰好配送一次：
\[
\sum_{p:b\in p}y_p=1.
\]
采用词典序目标：
\[
\boxed{
\min_{\rm lex}(N,E,T^{sum})
}.
\]
即首先最少架次，其次在最少架次条件下降低总能耗，最后缩短累计作业时间。

\subsection{精确求解与交叉验证}

Q1分别建立原始箱级MILP、对称压缩MILP和剩余需求状态DP。原始箱级模型直接枚举原始箱号子集；压缩MILP消除同区同类等价箱号对称；DP以剩余需求向量递推。15个服务区的三种精确方法得到相同词典序目标，因此最终结果不是单一求解器的一次输出。

\begin{table}[H]
\centering
\caption{Q1三种精确方法的交叉验证}
\begin{tabular}{lll}
\toprule
方法 & 主要作用 & 结果\\
\midrule
原始箱级MILP & 原问题精确基准 & 与另外两种方法一致\\
对称压缩MILP & 主精确复现 & 与箱级MILP一致\\
动态规划DP & 独立状态递推与Pareto枚举 & 与MILP一致\\
\bottomrule
\end{tabular}
\end{table}

\subsection{18架次下界与正式主方案}

三类机型中额定载荷最大为80 kg。S001总需求质量154 kg，S002和S003均为81 kg，因此这三个服务区至少各需2架次；其余12个服务区至少各需1架次。于是
\[
N\ge2+2+2+12=18.
\]
精确求解找到满足全部质量、体积与能源条件的18架次方案，因此
\[
\boxed{N^*=18}.
\]

固定18架次后继续最小化总能耗和累计作业时间，得到
\[
\boxed{
N=18,\quad
E=59.131296\ {\rm kWh},\quad
T^{sum}=546.266929\ {\rm min}
}.
\]
机型构成为
\[
A/B/C=0/9/9.
\]

\subsection{强基线与多目标权衡}

为了验证达到18架次后机型和箱组仍具有优化价值，正文优先与强基线比较。全C型精确方案同样为18架次，但能耗约为73.934908 kWh；每个服务区限定单一机型的精确策略约为63.068749 kWh；较强FFD类基线约为64.250537 kWh，而逐架次混合机型精确方案为59.131296 kWh。因此，相同架次数下，精确混合选型仍有明显节能价值。

完整不同Pareto目标向量只有两个：
\[
(18,\ 59.131296,\ 546.266929)
\]
和
\[
(19,\ 59.033942,\ 573.999623).
\]
增加1个架次可节省约0.097354 kWh，但累计作业时间增加约27.733 min，因此
\[
\boxed{\text{最少架次}\neq\text{最低能耗}}.
\]

\figplaceholder{图4-3占位：18架次主方案与19架次节能方案的离散Pareto权衡。}

\subsection{安全余量敏感性}

对所有满足额定质量与体积条件的候选批型定义
\[
\rho_p=1-\frac{E_p}{E_m},
\]
则统一安全余量 $\alpha$ 下有
\[
p\text{可行}\iff\alpha\le\rho_p.
\]
完整分析得到569个不同候选事件，其中8个事件会改变全局最少架次数。

\begin{table}[H]
\centering
\caption{统一安全余量引起的主要架次数临界事件}
\begin{tabular}{lcrr}
\toprule
临界余量/\% & 服务区 & 阈值处架次 & 严格右侧\\
\midrule
23.088350 & S004 & 18 & 19\\
27.049854 & S008 & 19 & 20\\
33.625809 & S008 & 20 & 21\\
33.896918 & S003 & 21 & 22\\
34.373532 & S004 & 22 & 23\\
34.391873 & S002 & 23 & 24\\
34.777897 & S008 & 24 & 25\\
35.267807 & S008 & 25 & 不可行\\
\bottomrule
\end{tabular}
\end{table}

当安全余量严格超过约35.2678\%时，S008的一只14 kg饮用水箱即使单独配送，也无法由任一机型完成安全往返；由于货箱不可拆，整个任务直接失去可行解。由此可见，连续安全要求会通过不可拆箱约束转化为离散架次数跃迁。

\figplaceholder[6.0cm]{图4-4占位：安全余量--最少架次数阶梯图及临界事件轨。}

% ============================================================
\section{问题二：异构多架次运输与共享电池联合调度}
\label{sec:q2}

\subsection{从单架次可行到系统级可执行}

Q1中一个运输任务只要自身满足物理约束即可视为可行；Q2中任务需要由有限实体无人机和共享电池真实执行。飞机返航并不意味着原电池已经恢复，因此
\[
\boxed{
\text{单架次物理可行}
\not\Rightarrow
\text{完整时间表可执行}
}.
\]

Q2需要同时决定货箱组批、访问路线、机型、实体飞机、电池以及任务开始时刻。

\figplaceholder{图5-1占位：Q2从运输任务构造到飞机--电池双资源联合排程及上下界认证的整体框架。}

\subsection{运输任务与双资源占用}

第3章统一任务表示仍写为
\[
r=(g_r,\mathcal B_r,\pi_r),
\]
并已计算完整作业时长 $p_r$、能耗 $e_r$、送达偏移 $\delta_{rb}$ 和返航SOC。

设任务在 $s_r$ 开始准备。实体无人机占用区间为
\[
I_r^U=[s_r,s_r+p_r),
\]
若返航后的充电恢复时间为 $c_r$，则电池占用区间为
\[
I_r^B=[s_r,s_r+p_r+c_r).
\]
因此
\[
\boxed{I_r^U\neq I_r^B}.
\]

\figplaceholder{图5-2占位：单个运输任务对实体无人机与共享电池的不同占用区间。}

\subsection{箱级时限与词典序目标}

若箱子 $b$ 由任务 $r$ 运输，则实际送达时刻为
\[
C_b=s_r+\delta_{rb}.
\]
迟到量与加权迟到为
\[
L_b=\max(0,C_b-d_b),
\qquad
L_w=\sum_b\omega_bL_b.
\]
Q2采用词典序目标：
\[
\boxed{
\min_{\rm lex}(L_w,T_{\max},E,N)
}.
\]
首先保证救援物资不迟到，再在零迟到最优面上最小化所有运输任务最后返航时间。

\subsection{联合调度模型与求解框架}

定义任务选择变量、实体飞机指派变量和电池指派变量，分别保证每个货箱恰好覆盖一次、每个已选任务由唯一兼容实体飞机和唯一兼容电池执行。使用同一飞机的任务区间 $I_r^U$ 不得重叠，使用同一电池的任务区间 $I_r^B$ 不得重叠。

完整组合空间同时具有组批、路线、机型和排程耦合，其资源占用与并行调度结构可由经典调度框架刻画\cite{pinedo2016scheduling}。因此采用“高质量可行上界+全局有效下界”的双路线：上界侧通过问题特定的破坏--修复、完整物理重建和事件排程寻找真实可执行方案；下界侧通过任务列松弛、完整定价和互补分支构造证书。大规模列式整数规划中的列生成与branch-and-price思想参见\cite{barnhart1998branchprice}。

\figplaceholder{图5-3占位：Q2高质量上界搜索、精确事件排程与下界认证的双路线框架。}

\subsection{23架次正式主方案}

第一层加权迟到非负，而正式方案达到
\[
L_w=0,
\]
因此
\[
\boxed{L_w^*=0}.
\]

当前正式主方案包含23个运输架次，
\[
A/B/C=11/6/6,
\]
使用8架实体运输无人机和14组共享电池。最后返航时间为
\[
\boxed{
T_{\max}=5693.231489\ {\rm s}
=94.887191\ {\rm min}
},
\]
总运输能耗为
\[
\boxed{
E=66.214460\ {\rm kWh}
}.
\]
最紧硬截止余量约为202.096 s，最低返航SOC为20.097383\%。

\begin{table}[H]
\centering
\caption{Q2正式主方案主要指标}
\begin{tabular}{lr}
\toprule
指标 & 结果\\
\midrule
加权迟到 & 0\\
运输架次 & 23\\
A/B/C架次 & 11/6/6\\
最后返航时间/s & 5693.231489\\
总运输能耗/kWh & 66.214460\\
实体运输机 & 8\\
共享电池 & 14\\
最紧硬截止余量/s & 202.096058\\
最低返航SOC & 20.097383\%\\
\bottomrule
\end{tabular}
\end{table}

\figplaceholder[7.0cm]{图5-5占位：8架实体运输无人机与14组共享电池的联合事件甘特图。}

\subsection{固定B型任务结构下的离散瓶颈}

当前6个B型任务完整作业时间总和为
\[
11133.182841\ {\rm s}.
\]
2架B型飞机的平均负载下界为
\[
5566.591421\ {\rm s},
\]
但6个不可拆任务不能连续均分。对固定B型任务进行精确两机分配，得到更强下界
\[
\boxed{
5693.231489\ {\rm s}
},
\]
而当前时间表恰好达到该值。因此，在保持当前6个B型任务的箱组、机型和路线不变时，仅重新排序已经不能继续缩短当前工期。进一步改善必须改变跨任务货箱、机型或访问结构。该结论仅适用于固定B型任务结构，并非原Q2的全局makespan下界。

\figplaceholder{图5-9占位：固定B型任务结构下6个任务在2架B型无人机上的离散负载瓶颈。}

\subsection{22架次可行见证与目标冲突}

进一步构造出一份真实可执行22架次方案：
\[
A/B/C=11/4/7,
\]
\[
T_{\max}=7767.058323\ {\rm s},
\qquad
E=66.323024\ {\rm kWh},
\qquad
L_w=0.
\]
因此23架次不是最少可行架次数。但当前22架次见证的工期比23架次主方案长约34.564 min，说明
\[
\boxed{
\text{更少架次}
\not\Rightarrow
\text{更短系统工期}
}.
\]
这里的7767.058 s只是一个可行见证，不是所有22架次方案的下界。

\figplaceholder{图5-6占位：22架次可行见证与23架次主方案的工期--能耗离散对照。}

\subsection{全局有效下界与认证区间}

仅搜索不到更好方案不能证明当前方案接近全局最优。为此，将完整运输任务作为列构造set-partitioning松弛，并通过完整定价和互补分支覆盖得到全局有效下界。

当前正式Q2证据版本为 Q2-EVIDENCE-V7-FOCUSED。最终得到
\[
\boxed{
LB=5433.956428\ {\rm s}
}
\]
以及有理上界
\[
\boxed{
UB=5693.231490\ {\rm s}
}.
\]
因此
\[
\boxed{
5433.956428\le T^*\le5693.231490
}.
\]
区间宽度为
\[
259.275062\ {\rm s},
\]
以上界为分母的证据间隙为
\[
\boxed{
\frac{UB-LB}{UB}=4.5541\%
}.
\]

最终覆盖包含809个互补分支与810份叶证书，对全部叶证书重新进行2430次完整机型定价，审计失败为0。上下界尚未闭合，因此不能宣称5693.231489 s已证明为全局最优。

\figplaceholder{图5-8占位：Q2零加权迟到最优面上的确定性完成时间上下界认证区间。}

\subsection{补充搜索与计算证据}

围绕当前关键任务进行了330个定向局部邻域重构，未发现严格改善上界；另有一个限时模型保持未知，不将其解释为不可行。更大的跨资源链6/7任务有限邻域同样未产生工期改善。这些结果仅支持局部稳定性，不替代全局证明。

针对证明阶段的pricing kernel，在24个固定真实定价输入上，支配强化将各组中位耗时之和由6.323099 s降至2.026581 s，比值为3.1201；累计生成标签由5,706,982降至2,061,175，减少63.88\%，且24个输入的最小约化成本全部保持一致。随后完整重定价810份证书、2430次调用后原下界仍有效。该3.1201倍仅表示固定pricing输入的合计中位耗时比，不表示整个Q2求解器提速3.12倍。

% ============================================================
\section{问题三：连续通信约束下的运输--中继联合优化}
\label{sec:q3}

\subsection{从运输可执行到全过程通信可用}

Q2运输时间表仍隐含“运输无人机全过程通信可用”的条件。山区地形可能使无人机在爬升、巡航、下降或返回途中出现通信遮挡，因此
\[
\boxed{
\text{运输计划可执行}
\not\Rightarrow
\text{运输全过程通信可用}
}.
\]

设运输无人机、中继无人机和地面站分别为 $U,R,G$。直连可用时优先采用 $U\leftrightarrow G$；直连失效时允许一次中继
\[
U\leftrightarrow R\leftrightarrow G.
\]
中继无人机同样需要起飞、到位、服务、返航和能源恢复，因此Q3本质上是运输轨迹、连续通信和中继资源时序的联合问题。

\figplaceholder{图6-1占位：从运输可执行到连续通信保障的模型升级。}

\subsection{链路预算与连续通信状态}

无人机辅助无线通信具有高度可控的三维移动性，同时受到视距传播、能量与时空资源耦合约束\cite{zeng2016uavcom,zeng2019sky}。自由空间传播损耗计算采用标准链路预算口径\cite{iturp525}。三类链路允许损耗上限为
\[
L_{UG}^{\max}=122\ {\rm dB},\qquad
L_{UR}^{\max}=116\ {\rm dB},\qquad
L_{RG}^{\max}=126\ {\rm dB}.
\]
定义链路余量
\[
M_{ab}(t)=L_{ab}^{\max}-L_{ab}(t).
\]
若考虑额外统一传播损耗 $\Delta$，则链路需满足
\[
M_{ab}(t)\ge\Delta.
\]
通信可行同时要求链路预算满足且传播路径未被地形阻挡。

对运输任务 $r$，无人机位置随时间连续变化，记为 $\mathbf x_r(t)$。定义直连不可用时间集合
\[
\mathcal W_r
=
\{t\in\mathcal T_r:\chi_{UG}(t)=0\}
=
\bigcup_{\ell=1}^{n_r}[a_{r\ell},b_{r\ell}].
\]
每个区间对应一项真实中继服务需求。本文不以固定步长采样替代连续条件，而根据轨迹分段、传播几何和通信状态变化构造区间并检查临界端点。

\figplaceholder{图6-2占位：代表运输任务沿时间轴的直连状态与中继需求区间。}

\subsection{中继候选位置与真实中继任务}

候选悬停位置集合记为 $\mathcal P^R$，位置 $p$ 的三维坐标为 $\mathbf z_p$。若在某通信需求区间内同时满足 $U\leftrightarrow R$ 和 $R\leftrightarrow G$，则候选位置能够在几何层面覆盖该需求。

正式有限域下界使用5755个明确候选位置。覆盖可行只说明“如果中继已经在线则能够通信”，并不保证中继可以在正确时间到位。因此，一个真实中继任务还必须包含
\[
\text{准备}
\rightarrow
\text{飞往中继点}
\rightarrow
\text{悬停服务}
\rightarrow
\text{返航}
\rightarrow
\text{周转与能源恢复}.
\]

中继实体和能源组件分别具有独立占用区间，最终充电不计入联合完成时间，但计入资源占用和Q4资源核算。

\figplaceholder{图6-3占位：中继架次从出发、到位、保障、返航到能源恢复的完整生命周期。}

\subsection{运输--通信联合模型与求解闭环}

Q3同时决定运输任务结构、运输时刻、中继位置、中继服务时间、中继实体和能源组件。Q2方案是高质量起点，但不是不可改变的冻结输入；如果纯运输最优结构难以获得通信保障，允许重新调整部分箱组、路线、机型和时间。

定义联合完成时间
\[
T_{\max}^{joint}
=
\max\{T_{\max}^{U},T_{\max}^{R}\}.
\]
在维持零加权迟到和硬截止可行的条件下，主方案首先最小化联合最后返航时间。通信鲁棒性通过额外损耗 $\Delta$ 构造独立情景，而不与时间、能耗强行加权。

整体求解形成
\[
\boxed{
\text{运输计划}
\rightarrow
\text{通信需求}
\rightarrow
\text{中继候选}
\rightarrow
\text{联合排程}
\rightarrow
\text{连续验证}
\rightarrow
\text{反向调整}
}.
\]

\figplaceholder{图6-4占位：运输--通信联合求解与连续验证闭环。}

\subsection{正式运输--通信联合方案}

正式 \texttt{time} 方案包含
\[
\boxed{
23\text{个运输架次}
+
4\text{个中继架次}
}.
\]
运输部分使用8架实体运输无人机和14组运输电池；中继部分实际使用2架中继无人机和3组中继能源组件。

联合最后返航时间为
\[
\boxed{
T_{\max}^{joint}=6340.439339\ {\rm s}
=105.674\ {\rm min}
}.
\]
运输能耗为
\[
66.376821\ {\rm kWh},
\]
中继能耗为
\[
2.931710\ {\rm kWh},
\]
总能耗为
\[
\boxed{
69.308531\ {\rm kWh}
}.
\]
加权迟到与硬迟到均为0，最小硬截止余量为30 s，最低运输返航SOC为20.097383\%。

\begin{table}[H]
\centering
\caption{Q3正式联合方案主要指标}
\begin{tabular}{lr}
\toprule
指标 & 结果\\
\midrule
运输架次 & 23\\
中继架次 & 4\\
联合最后返航/s & 6340.439339\\
运输能耗/kWh & 66.376821\\
中继能耗/kWh & 2.931710\\
总能耗/kWh & 69.308531\\
加权迟到 & 0\\
最小硬截止余量/s & 30\\
最低运输返航SOC & 20.097383\%\\
\bottomrule
\end{tabular}
\end{table}

\figplaceholder[7.0cm]{图6-5占位：正式运输轨迹、中继悬停位置及实际通信保障关系。}

\subsection{关键资源恢复机制}

一次针对S001两项C型运输任务的货箱调整，使关键前驱任务载荷由64 kg下降至53 kg，电池充电恢复时间由
\[
1528.177480\ {\rm s}
\]
降至
\[
1471.152841\ {\rm s},
\]
关键电池提前
\[
\boxed{
57.024638\ {\rm s}
}
\]
恢复可用，最终使联合工期缩短约
\[
\boxed{
54.651611\ {\rm s}
}.
\]
但运输总能耗略有增加。因此
\[
\boxed{
\text{最低总能耗}
\not\Rightarrow
\text{关键资源最早恢复}
}
\]
且关键资源恢复时刻可能比总能耗更直接决定系统工期。

\figplaceholder{图6-6占位：关键运输任务减载--电池提前恢复--联合工期缩短的传导机制。}

\subsection{连续通信完整验证}

正式联合方案对真实运输时序重新进行区间、端点、轨迹阶段切换和通信提供方切换检查。主验证器完成19215项联合验证检查，独立过程再次生成36995个核验点，未发现正式方案中的通信中断。数量本身不是核心，关键是验证对象覆盖实际连续区间和临界端点，而非仅服务区节点或固定步长采样。

\figplaceholder{图6-7占位：代表运输任务在完整执行过程中的Direct/Relay通信状态切换。}

\subsection{有限候选域中的4中继架次下界}

固定正式运输几何航迹后，共整理得到196个真实通信保障需求。对5755个候选悬停位置计算集合覆盖关系。构造一组非负需求权重，使
\[
\sum_j\omega_j=3.5,
\]
且任意一个候选中继架次能够覆盖的这些权重总和均不超过1。因此3个中继架次最多覆盖总权重3，无法覆盖全部需求，故
\[
N_R\ge\lceil3.5\rceil=4.
\]
正式方案恰好使用4个中继架次，所以在
\[
\boxed{
\text{固定运输轨迹}
+
5755\text{候选位置}
}
\]
组成的有限候选域中达到中继架次数下界。该证明不能推广到自由连续三维选址。

\subsection{通信裕度与时间--能耗权衡}

正式最快方案对应 $\Delta=0$：
\[
T=6340.439339\ {\rm s},\qquad
E=69.308531\ {\rm kWh}.
\]
0.50 dB完整可行方案 \texttt{robust\_fast} 为
\[
T=6369.026365\ {\rm s},\qquad
E=69.314466\ {\rm kWh}.
\]
0.55 dB完整可行方案 \texttt{robust\_margin} 为
\[
T=6395.090950\ {\rm s},\qquad
E=69.235215\ {\rm kWh}.
\]

\begin{table}[H]
\centering
\caption{不同通信裕度完整可行方案}
\begin{tabular}{lrrr}
\toprule
方案 & 额外损耗/dB & 联合工期/s & 总能耗/kWh\\
\midrule
time & 0 & 6340.439339 & 69.308531\\
robust\_fast & 0.50 & 6369.026365 & 69.314466\\
robust\_margin & 0.55 & 6395.090950 & 69.235215\\
\bottomrule
\end{tabular}
\end{table}

因此通信裕度具有可量化时间代价，但能耗并不随裕度简单单调增加。对固定新四点布局和固定运输轨迹，进一步得到
\[
0.55\le\Delta^*\le0.555382021\ {\rm dB},
\]
该区间不是自由三维选址下的全局裕度界。

\figplaceholder{图6-8占位：0/0.50/0.55 dB三种完整可行方案的联合工期与总能耗离散权衡。}

\subsection{静态代理评分与动态可执行性的差异}

对北部中继区域进行关键点引导搜索时，静态代理评分的最好样本达到0.573414 dB，中位数达到0.555772 dB，高于均匀采样。但该评分假设中继始终在线；将高分候选放回真实中继飞行、服务窗口、实体资源和能源组件排程后，并没有得到新的正式完整可行方案。

因此
\[
\boxed{
\text{更高静态通信评分}
\not\Rightarrow
\text{更好的完整联合方案}
}.
\]
这一负结果直接说明Q3不能退化为单纯的中继位置覆盖优化。

% ============================================================
\section{问题四：独立任务分组下的异构资源配置与池化损失}
\label{sec:q4}

\subsection{从全局共享到独立执行}

Q3正式方案默认所有运输无人机、电池、中继无人机和能源组件属于同一个共享资源池。若服务区被划分为多个相互独立的执行单元，则不同组之间不能继续共享实体资源，原本跨服务区存在的错峰复用链会被组织边界切断。因此
\[
\boxed{
\text{全局共享条件下资源充足}
\not\Rightarrow
\text{独立分组后各组资源仍然充足}
}.
\]

Q4严格冻结Q3 \texttt{time} 方案的货箱组批、运输机型、访问顺序、绝对任务时刻、中继位置、服务时段和实际通信保障关系，只改变服务区归属。

\figplaceholder{图7-1占位：全局共享与独立分组条件下资源复用关系的变化。}

\subsection{不可拆任务块}

由于同一运输架次和同一逻辑中继任务不能跨组拆分，运输关联与实际中继保障关联共同形成三个不可拆任务块：
\[
B_1=\{S001\},
\]
\[
B_2=
\{S002,S003,S004,S005,S007,S008,S009,S010,S011,S012,S013,S014,S015\},
\]
\[
B_3=\{S006\}.
\]
因此严格Q4不是直接划分15个独立服务区，而是对3个不可拆任务块进行非空分组。

\figplaceholder{图7-2占位：Q3冻结运输与中继关系形成的三个不可拆服务区块。}

\subsection{组内最小分类型资源需求}

资源类型集合为
\[
\mathcal P=
\{U_A,U_B,U_C,B_A,B_B,B_C,U_R,K_R\},
\]
原始库存向量为
\[
\mathbf I=(4,2,2,6,4,4,2,6).
\]

对任务组 $g$ 和资源类型 $p$，Q3已确定固定占用区间 $I_{jp}=[s_{jp},r_{jp})$。定义时刻 $t$ 的同时占用数
\[
n_{gp}(t)
=
\sum_{j\in\mathcal J_g}
\mathbf 1\{t\in I_{jp}\}.
\]
由于同型资源可替换且任务时间已固定，资源区间可视为区间图；其最少着色数等于最大同时重叠数\cite{fulkerson1965interval}。因此组内类型 $p$ 的最小资源数量恰好等于最大同时重叠数：
\[
\boxed{
N_{gp}^{min}
=
\max_t n_{gp}(t)
}.
\]

\figplaceholder{图7-3占位：固定占用区间最大重叠数与最小资源配置关系。}

\subsection{异构资源缺口与评价顺序}

若共有 $K$ 个独立组，类型 $p$ 的总需求为
\[
D_p=\sum_{g=1}^{K}N_{gp}^{min},
\]
缺口为
\[
G_p=\max(0,D_p-I_p).
\]
形成异构缺口向量 $\mathbf G$ 和总缺口
\[
G_\Sigma=\sum_pG_p.
\]

正式方案首先最小化资源缺口，其次最小化总配置数量，再考虑工作量变异系数 $CV_W$：
\[
\boxed{
\text{资源缺口}
\rightarrow
\text{总配置}
\rightarrow
\text{工作量均衡}
}.
\]
由于不同类型资源不能互换，即使
\[
\sum_pD_p\le\sum_pI_p,
\]
也不能推出所有类型均无缺口。

\subsection{严格模型完整枚举与两组结果}

经过任务块压缩后，两组仅有3种无标签合法分区，三组只有1种，因此全部采用完整枚举，不需要启发式算法。

两组全部合法方案如下。

\begin{table}[H]
\centering
\caption{严格两组模型的全部合法分区}
\begin{tabular}{lrrr}
\toprule
分组结构 & 总配置 & 总缺口 & 工作量CV\\
\midrule
其余14区 / S006 & 29 & \textbf{2} & 0.9424\\
S001+S006 / 其余13区 & 33 & 6 & \textbf{0.7837}\\
S001 / 其余14区 & 32 & 5 & 0.8413\\
\bottomrule
\end{tabular}
\end{table}

按正式词典序，资源优先方案为
\[
\boxed{
\{\text{除S006外的14个服务区}\}
\mid
\{S006\}
}.
\]
其需求向量为
\[
\mathbf D^{(2)}=(4,2,3,6,4,5,2,3),
\]
相对库存缺口为
\[
\boxed{
\mathbf G^{(2)}=(0,0,1,0,0,1,0,0)
}.
\]
即缺1架C型运输无人机和1组C型运输电池，总缺口2。

\begin{table}[H]
\centering
\caption{全局共享、两组独立与三组独立的分类型资源需求}
\begin{tabular}{lrrrr}
\toprule
资源类型 & 原库存 & Q3全局共享 & 两组独立 & 三组独立\\
\midrule
A型运输机 & 4 & 4 & 4 & 5\\
B型运输机 & 2 & 2 & 2 & 2\\
C型运输机 & 2 & 2 & 3 & 5\\
A型电池 & 6 & 6 & 6 & 7\\
B型电池 & 4 & 4 & 4 & 4\\
C型电池 & 4 & 4 & 5 & 6\\
中继无人机 & 2 & 2 & 2 & 2\\
中继能源组件 & 6 & 3 & 3 & 3\\
\midrule
总配置 & 30 & 27 & 29 & 34\\
\bottomrule
\end{tabular}
\end{table}

两组总配置29小于库存总数30，但仍缺C型机和C型电池，说明
\[
\boxed{
\text{总资源数量充足}
\not\Rightarrow
\text{异构资源结构充足}
}.
\]

\figplaceholder{图7-4占位：两组独立条件下分类型资源需求与原库存比较。}

\subsection{资源复用链断裂机制}

Q3全局方案中，S006相关C型任务在1561.553868 s返航，随后S001相关任务可继续复用同一C型实体；分组后该跨组复用被切断，因此C型无人机需求从2增加为3。

类似地，一组C型电池在S006任务后充电，并在3098.508915 s左右恢复满电，后继C型任务几乎立即开始；分组后该跨组电池链也被切断，因此C型电池需求从4增加为5。

所以新增资源不是因为任务数变多，而是
\[
\boxed{
\text{原有资源复用边被组织边界切断}
}.
\]

\figplaceholder{图7-5占位：S006独立后C型无人机与C型电池复用链断裂机制。}

\subsection{资源最省与任务最均衡的冲突}

资源缺口最少的两组方案工作量CV为0.9424，而更均衡的方案CV为0.7837，但总配置增加到33，库存缺口增加到6。说明
\[
\boxed{
\text{最少资源缺口}
\not\Rightarrow
\text{最均衡任务划分}
}.
\]
原因是工作量均衡关注总工作量，而资源配置取决于每种资源的时间峰值，两者并不等价。

\subsection{三组独立执行与池化损失}

严格三组分区唯一：
\[
\boxed{
\{S001\}
\mid
\{\text{其余13区}\}
\mid
\{S006\}
}.
\]
其需求向量为
\[
\mathbf D^{(3)}=(5,2,5,7,4,6,2,3),
\]
缺口向量为
\[
\boxed{
\mathbf G^{(3)}=(1,0,3,1,0,2,0,0)
}.
\]
即缺1架A型运输机、3架C型运输机、1组A型电池和2组C型电池，总缺口7，总配置34。

独立执行的资源代价可由
\[
\boxed{
\sum_{g=1}^{K}\max_t n_{gp}(t)
\ge
\max_t\sum_{g=1}^{K}n_{gp}(t)
}
\]
解释。左侧是各独立组分别按峰值配置再相加，右侧是共享资源池下的全局峰值；差值体现了失去跨组峰值错位共享所带来的资源池化损失。

\figplaceholder{图7-7占位：全局共享、两组独立与三组独立条件下的分类型资源需求变化。}

\subsection{充电占用反例}

若错误地认为电池在无人机返航后立即释放，而忽略充电恢复，则两组和三组总配置会被错误低估至约22和28；正确模型分别为29和34。由此可见
\[
\boxed{
\text{无人机返航}
\neq
\text{电池资源释放}
}.
\]
充电占用并非实现细节，而会直接改变组织配置结论。

\subsection{中继任务复制反事实}

正式主模型不允许原逻辑中继任务跨组复制。作为组织策略敏感性，若允许为每个涉及组完整复制中继任务，则分组自由度显著增加：两组有511种分区，三组有9330种分区，仍可完整枚举。

复制情景下三组最小资源缺口可由7降至4，总配置降至32，但需要新增一趟中继任务，并增加约0.8981 kWh中继能耗；全部复制情景中仍不存在依靠原库存实现零缺口的三组方案。因此复制不是免费改善，只是改变组织规则后的反事实方案。

总资源增加还应拆分为“任务复制成本”和“池化损失”。代表性复制情景中，原任务共享资源需求总数为27；加入复制任务后即使仍共享，需求已增至29；再禁止跨组共享后增至32，因此
\[
32-27
=
\underbrace{(29-27)}_{\text{复制任务成本 }2}
+
\underbrace{(32-29)}_{\text{池化损失 }3}.
\]

\subsection{库存敏感性与交叉验证}

对各类资源分别进行“任务开始前库存减少1个”的敏感性分析。运输无人机、运输电池或中继无人机减少时会增加缺口；中继能源组件因库存冗余，减少1组仍不改变正式主方案缺口。该分析讨论的是开工前可用库存，而不是执行过程中随机故障。

资源最小数量采用事件扫描、直接区间重叠、最早释放优先实体构造、兼容DAG最大匹配和独立MILP交叉核验，所得结果一致。详细全分区、全敏感性情景和实现审计移至附录。

% ============================================================
\section*{当前证据放置说明}
\addcontentsline{toc}{section}{当前证据放置说明}

为避免正文重复，公共AEQD/DEM/能耗验证统一放在第3章；Q1只保留精确求解一致性、18架次下界、强基线、Pareto与余量事件；Q2只保留完整资源时间表、固定B任务瓶颈、22架次见证和V7上下界；Q3保留连续通信验证、有限候选域4中继下界和鲁棒性；Q4保留完整枚举、资源向量、共享链断裂和复制/库存敏感性。完整日志、全量表、故障注入、证书明细和全分区应在后续附录中展开。


% ============================================================
\section{综合讨论、模型检验与决策启示}
\label{sec:discussion}

前四个问题分别从单架次安全运输、异构资源调度、连续通信保障和独立资源配置四个层面研究山区洪涝救援中的无人机协同任务。若仅分别观察各问结果，可以得到18架次、5693.231 s、4个中继架次以及2/7个资源缺口等具体数值；但将四个问题放在同一系统中考察后，可以发现更值得关注的是这些数值背后反复出现的结构规律。

本章不再重新建立模型或重复验证过程，而是从四个问题之间的共同机制出发，分析局部能力与系统执行、能耗与工期、静态指标与动态可执行性以及资源总量与资源结构之间的关系，并进一步讨论现有结论的证据强度、适用范围和实际救援决策意义。

\subsection{从局部可行到系统级可执行}

四个问题表面上分别关注组批、调度、通信和分组配置，但其共同主线可以概括为：
\[
\boxed{
\text{局部层面的可行性不会自动传递为更高层次的系统可执行性}
}.
\]

问题一首先判断某一批货物能否由某种机型安全完成单点往返。如果只考虑这一层，只需满足质量、体积和返航能源约束即可。

然而进入问题二以后，即使所有单架次任务分别物理可行，也仍然可能因为实体无人机不足、共享电池尚未恢复或任务时限冲突而无法构成完整时间表。因此
\[
\boxed{
\text{单架次安全可行}
\not\Rightarrow
\text{多任务联合可执行}
}.
\]

问题三进一步表明，一份运输时间表即使能够真实排入无人机和电池，也不能保证运输无人机在执行过程中始终具有通信链路。山区地形可能使无人机在航迹中间短暂失联，因此
\[
\boxed{
\text{运输排程可执行}
\not\Rightarrow
\text{全过程通信可执行}
}.
\]

而在问题四中，即使第三问已经形成完整的运输--通信联合方案，只要组织方式由一个共享资源池变为多个相互独立的任务组，原本能够跨服务区复用的无人机、电池和中继资源就可能无法继续共享。于是
\[
\boxed{
\text{全局共享可执行}
\not\Rightarrow
\text{独立分组后仍可执行}
}.
\]

因此，四问真正形成的是一条逐层收紧的可行性链：
\[
\boxed{
\text{单架次物理可行}
\rightarrow
\text{资源时间可行}
\rightarrow
\text{通信全过程可行}
\rightarrow
\text{组织独立可行}
}.
\]

每向后推进一层，都不是简单增加一个约束，而是重新改变了“什么才算一个真正可执行方案”的定义。

\figplaceholder{图8-1占位：四个问题中系统可行性的逐层收紧关系。}

\subsection{连续安全参数为何产生离散组批跃迁}

问题一中，返航安全余量 $\alpha$ 是连续变化的参数。随着 $\alpha$ 增加，无人机可用于任务执行的能量预算
\[
(1-\alpha)E_m
\]
连续下降，相应最大安全载荷也连续减小。

如果运输货物可以任意拆分，那么系统响应通常也会相对平滑；但本题货箱不可拆分，因此连续的物理能力变化最终会作用到离散组合结构上，由此形成
\[
\boxed{
\text{连续物理参数}
\rightarrow
\text{离散决策变化}
}.
\]

完整敏感性分析中虽然存在569个候选安全余量事件，但真正改变全局最少架次数的事件只有8个。例如，当安全余量跨过约
\[
23.08835\%
\]
后，最少架次数由18增加至19；继续提高安全余量后，又会在若干离散阈值处逐步增加。

当安全余量严格超过约
\[
35.2678\%
\]
时，S008中的一只14 kg饮用水箱即使单独运输也无法由任一机型安全完成往返，因此整个任务由“需要更多架次”直接转变为“不可行”。

这说明在实际任务决策中，安全参数的价值并不只体现为“安全余量每增加1\%会多消耗多少电”，更重要的是判断是否跨过了改变运输组织结构的临界点。

\subsection{能耗、资源恢复与系统工期并不等价}

Q1--Q3共同表明，“能耗更低”与“系统更快完成”并不是同一个目标。

在问题一中，18架次主方案总能耗为
\[
59.131296\ {\rm kWh},
\]
而增加到19架次后，可以将能耗进一步降至
\[
59.033942\ {\rm kWh}.
\]
因此
\[
\boxed{
\text{最少架次}
\not\Rightarrow
\text{最低能耗}
}.
\]

问题二进一步构造出22架次真实可行方案，但其联合运输工期为
\[
7767.058323\ {\rm s},
\]
明显长于23架次正式方案的
\[
5693.231489\ {\rm s}.
\]
所以
\[
\boxed{
\text{更少架次}
\not\Rightarrow
\text{更短系统完成时间}
}.
\]

问题三则揭示了更深一层的原因。对S001两项C型运输任务重新交换部分货箱后，关键前驱任务载荷由
\[
64\ {\rm kg}
\]
下降至
\[
53\ {\rm kg}.
\]
相应电池恢复时间由
\[
1528.177480\ {\rm s}
\]
下降至
\[
1471.152841\ {\rm s},
\]
使关键电池提前
\[
57.024638\ {\rm s}
\]
恢复可用，最终使联合工期缩短约
\[
54.651611\ {\rm s}.
\]
但这一调整并没有同步降低总运输能耗。

这说明，在存在资源复用和任务接续关系时，决定系统工期的更直接因素往往不是总能耗，而是关键资源链上的释放时刻与恢复时刻。因此，一个系统级调度方案至少需要区分单任务消耗、资源恢复速度和系统最后完成时刻。

\figplaceholder{图8-2占位：单任务能耗、关键资源恢复时刻与系统工期之间的传导关系。}

\subsection{静态能力评价不能代替动态执行验证}

在四问中反复出现“静态指标看起来很好，但放入真实时间表后不一定成立”的现象。

问题一中的最大安全载荷只回答某种机型在某个服务区理论上能够安全运输多重的货物，它无法回答所有运输任务能否同时由有限的实体无人机和电池执行。因此从Q1到Q2需要由静态任务可行进一步升级到动态排程可行。

类似地，在Q3中，一个中继候选位置可能具有较大的静态通信裕度，但如果中继无人机不能在运输无人机进入失联区域之前到达，或者中继能源组件正在执行其他任务，那么该位置仍然不能形成可执行方案。

关键点引导搜索得到的静态代理评分最高达到
\[
0.573414\ {\rm dB},
\]
高于均匀候选生成方法，但将这些高分位置重新代入完整中继飞行和资源排程后，并没有得到新的正式完整可行方案。因此
\[
\boxed{
\text{静态通信评价更好}
\not\Rightarrow
\text{联合通信任务更可执行}
}.
\]

问题二和问题三实际上揭示了相同的系统规律：
\[
\boxed{
\text{局部能力指标不能替代真实资源时间轴}
}.
\]

\subsection{异构资源结构与资源池化价值}

问题四进一步说明，资源可行性不仅由“有多少资源”决定，还由“资源是什么类型、什么时候被需要”共同决定。

严格两组方案的总资源配置需求为
\[
29,
\]
小于原始库存总数
\[
30.
\]
如果只比较总数，似乎原库存足以支持两组独立执行。但实际分类型需求表明，仍然缺少1架C型运输无人机和1组C型电池。

因此
\[
\boxed{
\sum_pD_p\le\sum_pI_p
\not\Rightarrow
D_p\le I_p,\quad\forall p
}
\]
也就是说
\[
\boxed{
\text{资源总量}
\neq
\text{资源结构}
}.
\]

另一方面，分组还会损失原本存在于全局资源池中的时间错峰复用能力。对于资源类型 $p$，共享资源池所需数量为
\[
N_p^{pool}
=
\max_t\sum_g n_{gp}(t),
\]
而独立执行所需数量为
\[
N_p^{ind}
=
\sum_g\max_t n_{gp}(t).
\]
由于
\[
\boxed{
\sum_g\max_t n_{gp}(t)
\ge
\max_t\sum_gn_{gp}(t)
}
\]
独立分组后的配置不会优于共享资源池。

因此，资源共享的价值并不仅仅是“大家可以用同一批设备”，其本质是允许某类资源在一个任务组的峰值结束以后继续服务另一个任务组。三组严格方案中，C型运输无人机需求由全局共享条件下的2架增加至5架，就是这一池化价值消失的典型表现。

\subsection{多目标决策中的不同“最优”}

四个问题还表明，随着决策层次改变，“最优”的含义也在改变。

\begin{table}[H]
\centering
\caption{不同问题中主要优化目标及其冲突关系}
\begin{tabular}{lll}
\toprule
问题 & 主要优先目标 & 主要冲突\\
\midrule
Q1 & 最少架次 & 架次数与能耗不等价\\
Q2 & 最短系统工期 & 架次数、能耗与工期不等价\\
Q3 & 最短联合工期 & 工期与通信裕度不等价\\
Q4 & 最少资源缺口 & 缺口与工作量均衡不等价\\
\bottomrule
\end{tabular}
\end{table}

这些指标并不存在天然统一的“最好”。因此本文没有构造类似
\[
aN+bE+cT+d\Delta
\]
的统一加权指标，因为不同量纲之间不存在自然、稳定的换算比例，任意权重反而会掩盖实际决策优先关系。采用词典序和独立备选方案，可以使每一层目标都保留明确解释。

\subsection{主要结论的证据等级}

不同章节采用的证明方式不同，因此需要明确区分不同结论的证据强度。

\begin{table}[H]
\centering
\caption{全文主要结论及证据范围}
\begin{tabular}{p{3.0cm}p{3.2cm}p{4.2cm}p{4.2cm}}
\toprule
结论 & 证据等级 & 可以支持的表述 & 不能扩大为\\
\midrule
Q1最少18架次 & 原Q1模型全局证明 & $N^*=18$ & 现实恶劣天气下绝对安全\\
Q1 Pareto结果 & 完整有限组合求解 & 当前模型下完整目标向量 & 任意外部物理模型均相同\\
Q2零加权迟到 & 全局证明 & $L_w^*=0$ & makespan全局最优\\
Q2工期 & 全局有效LB+完整可行UB & $5433.956428\le T^*\le5693.231490$ & 5693.231489已证明全局最优\\
Q3四中继架次 & 固定轨迹有限候选域证明 & 5755候选位置下至少需要4架次 & 任意连续三维位置至少需要4架次\\
Q3 0.55 dB & 完整可行见证 & 0.55 dB附加损耗可执行 & 自由选址全局最大裕度\\
Q4两/三组结果 & 严格模型完整枚举 & 缺口分别为2和7 & 重新设计Q3任务后仍必然为2和7\\
\bottomrule
\end{tabular}
\end{table}

其中Q2正式证据为
\[
\boxed{
5433.956428
\le
T^*
\le
5693.231490\ {\rm s}
}
\]
以上界为分母的认证间隙为
\[
\boxed{4.5541\%}.
\]
上下界尚未闭合，因此当前5693.231489 s方案应被表述为具有全局有效近优性证据的当前最优已知可行方案，而不能表述为已证明全局最优。

同理，Q3的4中继架次证明只在固定运输轨迹和5755个候选悬停位置的有限域内成立。明确这些边界能够避免把启发式搜索、有限域证明和原问题全局最优混为一谈。

\subsection{模型适用边界}

\subsubsection{物理环境边界}

本文采用题目给定的无人机性能、能耗参数、DEM和通信模型，并未额外引入实时风场、强降水空气动力影响、温度变化、电池健康状态衰减、GPS定位误差或随机机械故障。因此本文得到的是题目给定物理参数下的确定性任务规划结果，而不是现实环境下的飞行成功概率或航空安全认证。

\subsubsection{地形分辨率边界}

本文通过完整像元相交避免了稀疏DEM采样可能产生的漏检问题，但这种完整遍历只能保证已有DEM信息被完整使用，不能突破原始DEM本身的空间分辨率。因此
\[
\boxed{
\text{完整利用已有地形}
\neq
\text{获得无限精细地形}
}.
\]

\subsubsection{通信候选域边界}

Q3有限域下界建立在5755个明确中继候选位置和固定运输轨迹基础上。若允许中继在连续三维空间自由选址，可能存在当前候选集合中尚未包含的新位置，因此4中继架次不能被解释为连续三维自由选址条件下的全局下界。

\subsubsection{组织配置边界}

Q4严格冻结Q3正式运输和通信计划。因此Q4回答的是“已有救援预案被拆分成多个独立执行单元后，需要增加多少资源”，而不是“如果事先知道最终必须独立分组，并重新优化所有运输和中继任务，理论上最少需要多少资源”。

\subsection{参数敏感性与救援决策启示}

前文敏感性结果可以进一步转换为较直接的应急决策建议。

首先，安全余量提高会增强返航冗余，但其成本并不是平滑增加。管理者应优先判断是否接近会增加架次的结构临界点，而不仅比较单位安全余量对应的能耗变化。

其次，通信鲁棒性具有明确时间成本。正式最快方案在$\Delta=0$时联合工期为6340.439339 s；若要求额外0.50 dB统一传播损耗裕度，则工期增加约28.587 s；提高到0.55 dB时增加约54.652 s。因此可以根据灾区通信环境在任务完成时间与通信安全裕度之间进行透明权衡。

最后，资源补充应优先面向瓶颈类型。Q4两组独立执行时优先缺少C型运输机和C型电池，而中继能源组件仍有冗余。因此若只能有限增加装备，应优先补充瓶颈型号，而不是简单增加设备总数。

\subsection{模型特点与可扩展方向}

本文模型的主要特点不在于采用某一种特定优化算法，而在于保持四问之间的一致性。

首先，Q1--Q4共享同一套空间、地形、载荷、飞行时间和能耗模型，避免了不同问题使用不同物理口径。

其次，所有候选方案最终均需要回到真实时间轴验证资源可执行性，而不是停留在静态载荷、路径或通信覆盖层面。

再次，对于不同问题采用与其结构匹配的证据形式：小规模问题使用精确求解和完整枚举，大规模问题同时构造可行上界和有效下界，有限候选空间则明确限定证明范围。

若进一步面向现实救援应用，模型可以扩展至实时风场与天气情景、电池健康状态、多基地协同、执行过程在线重调度、随机设备失效、连续三维中继自由选址以及通信与运输联合鲁棒优化。

\subsection{本章小结}

综合四个问题，可以得到全文最重要的系统性认识：
\[
\boxed{
\text{局部可行性不会自动传递为系统级可执行性}
}.
\]

单架次安全能力需要经过不可拆组批后才能形成任务；运输任务还需要进入实体飞机和电池时间表；运输时间表仍需满足连续通信；全局联合方案在独立组织条件下又会因为资源共享链被切断而产生新的配置缺口。

同时，四问共同说明：
\[
\boxed{
\text{最低能耗、最少架次、最短工期、最大通信裕度和最均衡资源配置并不是同一目标}
}.
\]

系统决策的关键不是孤立优化某一个指标，而是识别当前决策层次中真正决定可执行性的瓶颈变量。

% ============================================================
\section{结论}
\label{sec:conclusion}

\subsection{总体结论}

本文针对山区洪涝灾害条件下无人机应急物资运输与通信保障问题，在统一地形、飞行、载荷和能源模型的基础上，依次研究了单点安全运输能力、多架次异构资源调度、连续通信中继以及独立任务组资源配置问题。

四个问题由低到高分别对应
\[
\boxed{
\text{安全能力}
\rightarrow
\text{时间协同}
\rightarrow
\text{通信连续}
\rightarrow
\text{资源独立}
}.
\]

研究结果表明，无人机救援系统的核心并不是单纯提高某一架无人机的性能，而是保证运输任务、能源恢复、通信保障和组织资源之间能够在同一时间轴上协调执行。

\subsection{问题一主要结论}

在20\%返航安全余量下，对15个服务区和A/B/C三类运输机共45个“服务区--机型”组合进行完整安全载荷计算，并在此基础上建立不可拆货箱精确组批模型。

三种独立精确方法得到一致结果，最少运输架次数满足
\[
\boxed{N^*=18}
\]
且达到严格下界。

在18架次条件下进一步优化能耗和累计作业时间，得到
\[
\boxed{E=59.131296\ {\rm kWh}}
\]
和
\[
\boxed{T^{sum}=546.266929\ {\rm min}},
\]
机型构成为
\[
A/B/C=0/9/9.
\]

同时存在另一非支配目标向量
\[
(19,\ 59.033942,\ 573.999623),
\]
说明
\[
\boxed{
\text{最少架次}
\not\Rightarrow
\text{最低能耗}
}.
\]

安全余量分析进一步得到8个改变全局最少架次数的关键阈值。当安全余量严格超过约35.2678\%时，任务整体失去可行性。因此，山区运输能力不能简单用无人机额定载荷表示，安全余量和离散货箱结构会共同决定实际所需架次数。

\subsection{问题二主要结论}

在问题一物理能力模型的基础上，进一步考虑实体无人机、共享电池、分段充电和箱级送达时限，构建异构多任务联合调度模型。

最终得到一份零加权迟到的23架次正式运输方案：
\[
\boxed{A/B/C=11/6/6}
\]
使用8架实体运输无人机和14组共享电池。

最后返航时间为
\[
\boxed{5693.231489\ {\rm s}}
\]
即约94.887 min，总运输能耗为
\[
\boxed{66.214460\ {\rm kWh}}.
\]

由于加权迟到非负，而当前方案达到$L_w=0$，因此第一层目标的全局最优值为
\[
\boxed{L_w^*=0}.
\]

对于第二层系统完成时间，最新有效认证区间为
\[
\boxed{
5433.956428
\le
T^*
\le
5693.231490\ {\rm s}
}
\]
以上界为分母的认证间隙为
\[
\boxed{4.5541\%}.
\]
由于上下界尚未闭合，因此本文不将5693.231489 s表述为已证明的全局最优工期。

另外构造了一份22架次真实可行见证：
\[
T=7767.058323\ {\rm s},
\qquad
E=66.323024\ {\rm kWh},
\]
说明23架次并非最少可行架次数，同时也说明减少架次数并不等价于缩短救援完成时间。

固定当前B型任务结构时，两架B型无人机的精确离散负载下界恰为5693.231489 s，表明当前任务结构中的瓶颈来自不可拆异构任务的离散负载，而不仅是总工作量。

\subsection{问题三主要结论}

在Q2运输调度基础上进一步加入地形遮挡、连续通信链路、中继悬停位置、中继无人机和能源组件，建立运输--通信联合优化模型。

正式 \texttt{time} 方案由
\[
\boxed{
23\text{个运输架次}
+
4\text{个中继架次}
}
\]
组成。

联合最后返航时间为
\[
\boxed{6340.439339\ {\rm s}}
\]
即约105.674 min，总能耗为
\[
\boxed{69.308531\ {\rm kWh}}.
\]
全部80个货箱继续保持零加权迟到。

通过对实际飞行轨迹进行区间和临界端点检查，正式方案全过程通信验证通过。

对于固定正式运输几何和5755个候选中继位置，共得到196个真实通信保障需求，并构造总权重3.5的对偶证书，从而证明
\[
\boxed{N_R\ge4}.
\]
而正式方案恰使用4个中继架次。因此，在固定运输轨迹和5755候选位置构成的有限域中，4个中继架次达到认证下界。该结论不推广到自由连续三维选址。

通信鲁棒性分析表明，$\Delta=0.50$ dB时存在完整可行方案
\[
T=6369.026365\ {\rm s},
\]
而$\Delta=0.55$ dB时仍存在完整可行方案
\[
T=6395.090950\ {\rm s}.
\]
这说明通信鲁棒性可以通过有限的额外时间成本获得。

此外，通过调整S001关键前驱任务，使关键电池提前约57.025 s恢复，最终联合工期缩短约54.652 s。因此关键资源恢复时刻比总能耗更能够直接解释联合系统的工期瓶颈。

\subsection{问题四主要结论}

冻结问题三正式运输任务、绝对时间和实际通信保障关系后，由运输关联和中继保障关系形成3个不可拆任务块。

在严格两组情形下，共有3种合法无标签分区。完整枚举得到资源缺口最小方案
\[
\boxed{
\{\text{除S006外14个服务区}\}
\mid
\{S006\}
}.
\]
总配置需求为29，相对原库存缺口为
\[
\boxed{
\mathbf G^{(2)}
=
(0,0,1,0,0,1,0,0)
}
\]
即需要新增1架C型运输无人机和1组C型电池。

严格三组合法分区唯一：
\[
\boxed{
\{S001\}
\mid
\{\text{其余13区}\}
\mid
\{S006\}
}.
\]
总资源配置为34，资源缺口为
\[
\boxed{
\mathbf G^{(3)}
=
(1,0,3,1,0,2,0,0)
}
\]
即需要新增1架A型运输机、3架C型运输机、1组A型电池以及2组C型电池，总缺口为7。

特别地，两组总配置只有29，而原库存总数为30，但系统仍然无法依靠原库存独立执行。因此
\[
\boxed{
\text{资源总量充足}
\not\Rightarrow
\text{异构资源结构充足}
}.
\]

允许完整复制中继任务作为反事实敏感性后，三组最低缺口可以从7下降至4，但会新增一趟中继任务及约0.8981 kWh中继能耗，并且仍然不存在原库存零缺口方案。由此说明，组织独立性的成本主要来自原有跨任务组资源复用链的断裂，而不是运输任务本身数量增加。

\subsection{跨问题主要研究发现}

综合四个问题，可以提炼出三条具有共同性的研究发现。

第一，连续物理约束可以产生离散系统响应。安全余量本身连续变化，但由于货箱不可拆和架次数为整数，最终运输组织会在若干阈值处发生离散跳变。因此
\[
\boxed{
\text{连续参数变化}
\not\Rightarrow
\text{系统决策连续变化}
}.
\]

第二，系统完成时间由关键资源链而不是单一局部指标决定。无论是Q2中的B型离散负载，还是Q3中的关键电池恢复，都表明系统工期受到资源接续关系控制。因此
\[
\boxed{
\text{最低局部能耗}
\not\Rightarrow
\text{最短系统工期}
}.
\]
真正需要关注的是关键资源何时释放并重新可用。

第三，异构资源的共享价值来自类型匹配和时间错峰。Q4中资源池化收益来源于不同任务峰值发生在不同时间，使同一实体能够跨任务连续复用。因此
\[
\boxed{
\text{资源池价值}
\neq
\text{简单资源数量之和}
}.
\]

\subsection{救援决策启示}

从实际救援组织角度，本文结果可以转化为以下总体原则。

在任务制定阶段，不应直接用额定载荷替代山区条件下的真实安全能力，而应先结合任务距离、地形和返航安全余量计算可执行载荷。

在运输调度阶段，不应仅追求减少架次数或降低总能耗，而应识别决定最后完成时间的关键无人机和电池恢复链。

在通信保障阶段，不应仅依据静态覆盖图选择中继位置，而应检查中继无人机能否在真实运输时间窗内到位并持续提供服务。

在独立资源配置阶段，不应根据设备总数判断库存是否充足，而应根据分类型资源峰值和原有资源复用链确定真正的短缺类型。

因此，山区无人机救援系统的关键不是单独提高某一架无人机的能力，而是保持
\[
\boxed{
\text{运输}
+
\text{能源}
+
\text{通信}
+
\text{组织资源}
}
\]
在统一时间轴上的协同关系。

最终，本文建立的整体逻辑可以概括为
\[
\boxed{
\text{安全能力}
\rightarrow
\text{时间协同}
\rightarrow
\text{通信连续}
\rightarrow
\text{资源独立}
}.
\]

只有将局部任务能力重新放回完整系统中评价，才能形成真正具有执行意义的山区洪涝无人机救援方案。


\clearpage
\renewcommand{\refname}{参考文献}
\bibliographystyle{unsrt}
\bibliography{references}

\clearpage
\appendix

\section{原始数据与公共物理模型核验}
\label{app:common}

\subsection{数据对象、单位与复现口径}

全文公共模型使用同一套节点、货箱、机型、DEM与任务参数。内部统一采用 m、s、kg、m$^3$、kWh 和 dB；正文显示时允许换算为 min，但后续计算始终使用未舍入原值。Q1公共物理验证不读取Q2--Q4求解结果，Q2--Q4均继承同一地形、航段与能耗口径。

\begin{table}[H]
\centering
\caption{公共模型关键核验结果}
\label{tab:app_common_validation}
\begin{tabular}{p{4.2cm}p{7.0cm}p{3.2cm}}
\toprule
核验对象 & 独立核验方式 & 结果\\
\midrule
TIF/MAT高程矩阵 & 逐像元比对 & 完全一致\\
AEQD水平距离 & GeographicLib独立测地线计算 & 最大差 $1.412\times10^{-9}$ m\\
DEM航段像元集合 & 网格边界交点遍历与条带区间求交 & 完全一致\\
45个安全载荷根 & Brent法与独立二分法 & 最大差 $1.843\times10^{-11}$ kg\\
Q1精确优化 & 原始箱级MILP、压缩MILP、DP & 词典序目标一致\\
自动化测试 & 42项测试 & 0失败、0错误、0跳过\\
\bottomrule
\end{tabular}
\end{table}

\subsection{公共物理错误注入}

为验证公共物理层检查器确实能够识别错误，而非仅重放既有结果，对以下七类错误进行了故障注入。

\begin{table}[H]
\centering
\caption{公共物理模型故障注入}
\label{tab:app_fault_injection}
\begin{tabular}{p{4.4cm}p{9.2cm}}
\toprule
错误类型 & 检出机制\\
\midrule
将载荷--航程指数 $3/2$ 错改为1 & 独立航程中间值与能耗重算不一致\\
遗漏返程爬升 & 往返能耗独立检查失败\\
返程错误沿用起飞载荷 & 逐段载荷重算失败\\
跨服务区混箱 & 服务区归属检查失败\\
重复分配同一货箱 & 原始货箱ID唯一计数失败\\
只检查质量、忽略体积 & 质量可行但体积超限的批型被拒绝\\
使用thin Bresenham替代完整像元遍历 & 航段相交像元集合不一致\\
\bottomrule
\end{tabular}
\end{table}

\subsection{复现性}

Q1在全新输出目录中从原始输入完整重跑，27个确定性数值文件与5幅生成图逐字节一致。求解耗时、平台元数据和可交换箱号代表不要求字节一致。公共模型的目标是保证“任务属性算得一致”，而各问最优性或近优性证据分别在后续附录中给出。

% ============================================================
\section{问题一完整能力、组批与敏感性结果}
\label{app:q1}

\subsection{45组最大安全载荷}

20\%返航安全余量下，15个服务区与3种机型的最大安全载荷如下。除表中低于额定载荷的6项外，其余组合均由额定载荷而非能源约束决定。

\begin{table}[H]
\centering
\caption{Q1各服务区三类机型最大安全载荷/kg}
\label{tab:app_q1_payload}
\begin{tabular}{lrrr}
\toprule
服务区 & A型 & B型 & C型\\
\midrule
S001 & 25.000000 & 30.000000 & 80.000000\\
S002 & 25.000000 & 30.000000 & 68.860343\\
S003 & 25.000000 & 30.000000 & 68.206491\\
S004 & 25.000000 & 30.000000 & 63.696955\\
S005 & 25.000000 & 30.000000 & 80.000000\\
S006 & 25.000000 & 30.000000 & 80.000000\\
S007 & 25.000000 & 30.000000 & 80.000000\\
S008 & 25.000000 & 28.801234 & 58.904361\\
S009 & 25.000000 & 30.000000 & 80.000000\\
S010 & 25.000000 & 30.000000 & 80.000000\\
S011 & 25.000000 & 30.000000 & 80.000000\\
S012 & 25.000000 & 30.000000 & 68.117193\\
S013 & 25.000000 & 30.000000 & 80.000000\\
S014 & 25.000000 & 30.000000 & 80.000000\\
S015 & 25.000000 & 30.000000 & 80.000000\\
\bottomrule
\end{tabular}
\end{table}

\subsection{18架次正式组批摘要}

完整箱号分配可由机器可读文件 \path{problems/D/q1/results/optimal_batches.csv} 复现。表\ref{tab:app_q1_batches}给出18个正式批次的主要物理属性。

\begin{longtable}{llllrrrr}
\caption{Q1正式18架次批次摘要}\label{tab:app_q1_batches}\\
\toprule
批次 & 区域 & 机型 & 质量/kg & 体积/m$^3$ & 能耗/kWh & 作业时间/s & 返航SOC/\%\\
\midrule
\endfirsthead
\toprule
批次 & 区域 & 机型 & 质量/kg & 体积/m$^3$ & 能耗/kWh & 作业时间/s & 返航SOC/\%\\
\midrule
\endhead
Q1-001 & S001 & C & 76 & 0.170 & 2.917582 & 1494.169 & 63.530\\
Q1-002 & S001 & C & 78 & 0.224 & 2.995769 & 1692.169 & 62.553\\
Q1-003 & S002 & B & 25 & 0.067 & 2.713422 & 1816.694 & 32.164\\
Q1-004 & S002 & C & 56 & 0.144 & 5.721225 & 2041.812 & 28.485\\
Q1-005 & S003 & B & 25 & 0.067 & 2.779917 & 2080.523 & 30.502\\
Q1-006 & S003 & C & 56 & 0.144 & 5.764965 & 2370.091 & 27.938\\
Q1-007 & S004 & C & 59 & 0.156 & 6.152932 & 2301.057 & 23.088\\
Q1-008 & S005 & C & 59 & 0.156 & 4.222193 & 1877.191 & 47.223\\
Q1-009 & S006 & C & 59 & 0.156 & 2.597988 & 1561.554 & 67.525\\
Q1-010 & S007 & C & 45 & 0.129 & 3.410894 & 1907.566 & 57.364\\
Q1-011 & S008 & C & 45 & 0.129 & 5.836012 & 2199.701 & 27.050\\
Q1-012 & S009 & B & 25 & 0.067 & 2.096234 & 1558.442 & 47.594\\
Q1-013 & S010 & B & 25 & 0.067 & 1.823196 & 1555.380 & 54.420\\
Q1-014 & S011 & B & 25 & 0.067 & 1.073952 & 1188.420 & 73.151\\
Q1-015 & S012 & B & 25 & 0.067 & 2.752414 & 1922.740 & 31.190\\
Q1-016 & S013 & B & 25 & 0.067 & 1.824759 & 1510.472 & 54.381\\
Q1-017 & S014 & B & 25 & 0.067 & 2.140388 & 1869.756 & 46.490\\
Q1-018 & S015 & B & 25 & 0.067 & 2.307453 & 1828.277 & 42.314\\
\bottomrule
\end{longtable}

\begin{figure}[H]
  \centering
  \includegraphics[width=0.82\textwidth]{map_single_routes.png}
  \caption{Q1基地O01至15个服务区的单点往返航线与DEM背景示意。该图仅用于空间关系展示，正式地形计算采用完整DEM像元相交。}
  \label{fig:app_q1_routes}
\end{figure}

\subsection{基线与Pareto结果}

\begin{table}[H]
\centering
\caption{Q1主要基线与正式方案}
\label{tab:app_q1_baseline}
\begin{tabular}{llrrr}
\toprule
机型策略 & 方法 & 架次 & 能耗/kWh & 累计作业时间/min\\
\midrule
A-only & exact & 52 & 112.180153 & 1488.763007\\
B-only & exact & 35 & 68.350657 & 927.625473\\
C-only & exact & 18 & 73.934908 & 563.979235\\
mixed & FFD-mass & 18 & 75.116697 & 563.979235\\
mixed & exact & 18 & \textbf{59.131296} & \textbf{546.266929}\\
\bottomrule
\end{tabular}
\end{table}

完整不同Pareto目标向量仅有两组：
\[
(18,\ 59.131296022053,\ 546.266929148593),
\]
\[
(19,\ 59.033942092203,\ 573.999623452513).
\]

\subsection{改变全局最少架次数的安全余量阈值}

\begin{table}[H]
\centering
\caption{Q1八个全局架次数临界阈值}
\label{tab:app_q1_thresholds}
\begin{tabular}{lcrr}
\toprule
安全余量阈值/\% & 触发区域 & 阈值处架次 & 严格右侧\\
\midrule
23.088350 & S004 & 18 & 19\\
27.049854 & S008 & 19 & 20\\
33.625809 & S008 & 20 & 21\\
33.896918 & S003 & 21 & 22\\
34.373532 & S004 & 22 & 23\\
34.391873 & S002 & 23 & 24\\
34.777897 & S008 & 24 & 25\\
35.267807 & S008 & 25 & 不可行\\
\bottomrule
\end{tabular}
\end{table}

完整569个候选事件与19次词典序最优代表变化由 \path{problems/D/q1/results/reserve_thresholds.csv} 和 \path{problems/D/q1/results/reserve_optimal_segments.csv} 保存。

% ============================================================
\section{问题二任务构造、调度结果与认证证据}
\label{app:q2}

\subsection{23个正式运输任务}

表\ref{tab:app_q2_tasks}直接对应 \path{q2_final/solution/main_23/decisions.json}。任务编号T01--T23仅用于论文展示，原始箱号保持不变。

\scriptsize
\begin{longtable}{cllp{8.2cm}}
\caption{Q2正式23个运输任务构成}\label{tab:app_q2_tasks}\\
\toprule
任务 & 机型 & 访问顺序 & 原始货箱ID\\
\midrule
\endfirsthead
\toprule
任务 & 机型 & 访问顺序 & 原始货箱ID\\
\midrule
\endhead
T01 & C & S001 & S001-MED-01, S001-MED-02, S001-WAT-06, S001-WAT-07, S001-WAT-08, S001-FOD-03, S001-HYG-01, S001-HYG-02\\
T02 & C & S001 & S001-WAT-02, S001-WAT-03, S001-WAT-04, S001-WAT-05, S001-FOD-02\\
T03 & B & S010 & S010-MED-01, S010-WAT-01, S010-FOD-01\\
T04 & B & S014 & S014-MED-01, S014-WAT-01, S014-FOD-01\\
T05 & B & S012 & S012-MED-01, S012-WAT-01, S012-FOD-01\\
T06 & A & S001 & S001-WAT-01, S001-FOD-01\\
T07 & A & S004 & S004-WAT-01, S004-FOD-01\\
T08 & A & S013 & S013-WAT-01, S013-FOD-01\\
T09 & A & S011$\rightarrow$S004 & S004-MED-01, S004-WAT-02, S011-MED-01\\
T10 & C & S006 & S006-MED-01, S006-WAT-01, S006-WAT-02, S006-WAT-03, S006-FOD-01, S006-HYG-01\\
T11 & C & S002 & S002-MED-01, S002-WAT-02, S002-WAT-03, S002-WAT-04, S002-FOD-01, S002-FOD-02, S002-HYG-01\\
T12 & C & S007$\rightarrow$S003 & S003-WAT-02, S003-WAT-03, S003-WAT-04, S003-FOD-01, S003-FOD-02, S003-HYG-01, S007-HYG-01\\
T13 & C & S005 & S005-MED-01, S005-WAT-01, S005-WAT-02, S005-WAT-03, S005-FOD-01, S005-HYG-01\\
T14 & A & S009$\rightarrow$S002$\rightarrow$S013 & S002-WAT-01, S009-MED-01, S013-MED-01\\
T15 & A & S007 & S007-MED-01, S007-WAT-01\\
T16 & A & S007 & S007-WAT-02, S007-FOD-01\\
T17 & B & S004 & S004-WAT-03, S004-HYG-01\\
T18 & B & S008 & S008-MED-01, S008-WAT-01, S008-FOD-01\\
T19 & B & S008 & S008-WAT-02, S008-HYG-01\\
T20 & A & S003$\rightarrow$S015 & S003-MED-01, S003-WAT-01, S015-MED-01\\
T21 & A & S015 & S015-WAT-01, S015-FOD-01\\
T22 & A & S009 & S009-WAT-01, S009-FOD-01\\
T23 & A & S011 & S011-WAT-01, S011-FOD-01\\
\bottomrule
\end{longtable}
\normalsize

\subsection{正式主方案与22架次见证}

\begin{table}[H]
\centering
\caption{Q2正式23架次方案与22架次可行见证}
\label{tab:app_q2_22_23}
\begin{tabular}{lrr}
\toprule
指标 & 正式23架次 & 22架次见证\\
\midrule
A/B/C架次 & 11/6/6 & 11/4/7\\
加权迟到/s & 0 & 0\\
makespan/s & 5693.231489 & 7767.058323\\
总能耗/kWh & 66.214460 & 66.323024\\
最紧硬截止余量/s & 202.096058 & 645.017134\\
最低返航SOC/\% & 20.097383 & 20.097383\\
\bottomrule
\end{tabular}
\end{table}

22架次方案只证明“22架次可行”，不证明22架次的最小makespan，也不能据此推出所有22架次方案均慢于23架次主方案。

\subsection{V7全局有效上下界}

\begin{table}[H]
\centering
\caption{Q2-EVIDENCE-V7-FOCUSED认证摘要}
\label{tab:app_q2_v7}
\begin{tabular}{lr}
\toprule
证据项 & 数值\\
\midrule
全局有效下界/s & 5433.956428\\
认证可行上界/s & 5693.231490\\
区间宽度/s & 259.275062\\
UB归一化间隙 & 4.554093092\%\\
互补分支数 & 809\\
叶证书数 & 810\\
完整机型定价调用 & 2430\\
审计失败数 & 0\\
原问题makespan全局最优 & 未证明\\
\bottomrule
\end{tabular}
\end{table}



\subsection{有限邻域与计算层负结果}

围绕当前主方案共检查330个定向局部邻域，329个有限模型返回不可行，1个达到时限后保持未知；不能将该结果解释为原问题局部或全局最优证明。pricing dominance实验只评价固定真实定价输入：24个输入的各组中位耗时之和由6.323099 s降至2.026581 s，比值为3.1201；生成标签由5,706,982降至2,061,175，最小约化成本保持一致。该3.1201倍不是整个Q2求解器速度比。

% ============================================================
\section{问题三运输--中继联合方案与验证}
\label{app:q3}

\subsection{当前正式release摘要}

当前正式版本为 \texttt{Q3-BOTTLENECK-V6/time}，物理模型版本为 \texttt{Q3-HOVER-COMPATIBLE-V5}。

\begin{table}[H]
\centering
\caption{Q3 V6正式主方案摘要}
\label{tab:app_q3_summary}
\begin{tabular}{lr}
\toprule
指标 & 数值\\
\midrule
运输架次 & 23\\
中继架次 & 4\\
实体运输无人机 & 8\\
运输电池 & 14\\
中继无人机 & 2\\
中继能源组件 & 3\\
联合最后返航/s & 5836.969929\\
运输能耗/kWh & 66.251280\\
中继能耗/kWh & 2.688136\\
总能耗/kWh & 68.939416\\
加权迟到 & 0\\
硬迟到箱数 & 0\\
最小硬截止余量/s & 360.866938\\
最低运输返航SOC/\% & 20.097383\\
\bottomrule
\end{tabular}
\end{table}

\subsection{中继转场高度口径}

中继转场采用
\[
H_R(p)=\max\left\{\max_{c\in\mathcal C(O01,p)}Z(c)+50,\ z_p,\ z_{O01}\right\}.
\]
该规则只作用于中继往返转场，普通运输航段仍使用沿线最高地形+50 m。旧release中将悬停高度额外截断在地形净空巡航面以下，因此其有限候选域证书不能自动继承到V6。

\subsection{连续通信验证}

V6主方案重新构造真实运输轨迹、完整中继需求区间与临界端点。主验证器完成17,756项物理、资源、区间和端点检查，失败0；独立时空点核验36,672点，未发现断链。保存证据包含700份完整区间证书和1,238份临界端点证书。

\subsection{关键保障切换}

关键C型运输任务在V5中从相对956.973974 s起必须依赖西侧中继；V6将此前段保障延长至1039.150288 s，同时使西侧中继提前5.821245 s就绪。两项共同解释约87.997559 s的同模型工期改善。该分析解释保存日程中的关键链，不构成原问题全局下界。

\subsection{时间--能耗与鲁棒备选}

\begin{table}[H]
\centering
\caption{Q3 V6已验证候选}
\label{tab:app_q3_robust}
\begin{tabular}{lrrr}
\toprule
方案 & 额外统一损耗/dB & 联合工期/s & 总能耗/kWh\\
\midrule
time & 0 & 5836.969929 & 68.939416\\
energy & 0 & 5911.548135 & 68.735791\\
robust025 & 0.25 & 5866.140527 & 69.031751\\
q4\_tradeoff & 0 & 5866.140527 & 68.975497\\
\bottomrule
\end{tabular}
\end{table}

旧+0.50/+0.55 dB结果属于历史release，不属于V6；当前只把+0.25 dB方案作为V6下完整核验的通信裕度备选。

% ============================================================
\section{问题四V6严格分区与资源向量}
\label{app:q4}

\subsection{严格模型不可拆块}

冻结Q3 V6/time任务、绝对时刻和实际通信保障关系后，严格模型形成
\[
B_1=\{S001\},\qquad
B_2=\{S002,S003,S004,S005,S007,S008,S009,S010,S011,S012,S013,S014,S015\},\qquad
B_3=\{S006\}.
\]
严格两组共有3种无标签分区，严格三组只有1种。

\subsection{当前正式资源向量}

\begin{table}[H]
\centering
\caption{Q4 V6严格分类型资源需求与库存}
\label{tab:app_q4_vectors}
\begin{tabular}{lrrrr}
\toprule
资源类型 & 原库存 & 全局共享 & 严格K2 & 严格K3\\
\midrule
A型运输机 & 4 & 4 & 4 & 5\\
B型运输机 & 2 & 2 & 2 & 2\\
C型运输机 & 2 & 2 & 3 & 5\\
A型电池 & 6 & 6 & 6 & 7\\
B型电池 & 4 & 4 & 4 & 4\\
C型电池 & 4 & 4 & 4 & 6\\
中继无人机 & 2 & 2 & 2 & 2\\
中继能源组件 & 6 & 3 & 3 & 3\\
\midrule
总配置 & 30 & 27 & 28 & 34\\
\bottomrule
\end{tabular}
\end{table}

严格两组的正式需求和缺口为
\[
\mathbf D^{(2)}=(4,2,3,6,4,4,2,3),\qquad
\boxed{\mathbf G^{(2)}=(0,0,1,0,0,0,0,0)}.
\]
严格三组的需求和缺口为
\[
\mathbf D^{(3)}=(5,2,5,7,4,6,2,3),\qquad
\boxed{\mathbf G^{(3)}=(1,0,3,1,0,2,0,0)}.
\]

两组正式分区为“除S006外14区 / S006”，总缺口1；三组唯一分区为“S001 / 其余13区 / S006”，总缺口7。

\subsection{跨问折中}

V6另保留一个 \texttt{q4\_tradeoff} 候选：联合工期5866.140527 s、总能耗68.975497 kWh，其严格三组缺口为6。它在Q3时间--能耗二维上不优于正式time方案，但说明Q3最短工期与Q4最少独立配置并不等价。

\subsection{历史敏感性边界}

旧Q4-SENSITIVITY-V2中的relay-copy、N-1、Pareto和库存敏感性使用旧Q3日程，未针对V6重新计算。它们仍保留在仓库历史目录，但不能与当前V6 strict资源向量混用。

\subsection{资源最小数量的交叉核验}

对固定任务组和固定资源类型，最小资源数量
\[
N_{gp}^{\min}=\max_t n_{gp}(t)
\]
由事件扫描和半开区间峰值计算得到，并在V6完整运行包中以实体资源链和独立最大匹配进行交叉核验。

% ============================================================

\section{复现、版本锁与证据清单}
\label{app:repro}

\subsection{正式版本入口}

\begin{table}[H]
\centering
\caption{四问正式结果入口}
\label{tab:app_release}
\begin{tabular}{lll}
\toprule
问题 & 正式版本/提交 & 论文复现入口\\
\midrule
Q1 & \texttt{f905acc8482966c...} & \path{problems/D/q1/}\\
Q2 & Q2-EVIDENCE-V7-FOCUSED & \path{problems/D/q2/evidence_v7/}\\
Q3 & Q3-BOTTLENECK-V6 / time & \path{problems/D/q3/release_v6/time/}\\
Q4 & Q4-STRICT-FROM-Q3-V6 & \path{problems/D/q4/results_q3_v6/}\\
\bottomrule
\end{tabular}
\end{table}

\subsection{哈希锁定的伴随归档}

\begin{longtable}{p{6.5cm}p{8.0cm}}
\caption{正式伴随归档SHA-256}\label{tab:app_archives}\\
\toprule
归档 & SHA-256\\
\midrule
\endfirsthead
\toprule
归档 & SHA-256\\
\midrule
\endhead
D题\_Q2定向证据强化V7\_源码证书与复核.zip &
\texttt{0804f989699dc506b3111dacc7408d9acdeb6e1115e7f32ec36934116de6ce58}\\
D题\_Q3正式封口\_FINAL\_源码结果与交接.zip &
\texttt{97f5ea8018a557d604bee87d6886675ecf7648f035e8df0f47f20ff6e7fbc56b}\\
D题\_Q4稳健性完善V2\_源码结果与验证.zip &
\texttt{9c8d1a6410c2de0a410318fc3525dc162613deabf4ba9dd1d59cdeb98437f5b4}\\
D题\_Q2外部方案复核\_代码明细与验证.zip &
\texttt{b41a0416313a123e734e01d7ca1672b99acee41f16edd79812302e1007610b77}\\
\bottomrule
\end{longtable}

\subsection{证据等级与最终提交前检查}

正式论文中只允许使用与证据强度相匹配的表述：
\begin{itemize}
  \item Q1最少18架次：原Q1模型下全局证明；
  \item Q2加权迟到0：第一层目标全局证明；
  \item Q2 makespan：有效全局下界与完整可行上界，未闭合；
  \item Q3 V6：连续区间/临界端点与独立重放证明保存方案完整可行；
  \item Q3固定结构：连续时间LP证明当前固定离散结构下的时序最优；
  \item Q3鲁棒性：仅$+0.25$ dB为当前V6完整可行见证，旧$+0.50/+0.55$ dB不继承；
  \item Q4 K2/K3：冻结Q3 V6/time日程后的strict完整分区枚举。
\end{itemize}

最终提交前应再次以 \path{paper_integration/FINAL_SELECTION.json} 为唯一选择入口，逐项核对正文、摘要、图表、附录和机器可读结果，避免旧版本数字残留。


\end{document}
